% Conjectura manuscript, CJ-0001: the Collins--Shor conjecture (both halves).
% DRAFT FOR INDEPENDENT REVIEW. NOT REVIEWED BY A HUMAN. NOT FOR DISTRIBUTION.
% Written on 2026-11-27 (session 14) by Claude, model claude-opus-5-5, from the
% project notes listed in Appendix A. The lower half is restated in a shorter,
% self-contained form (root weights only); the upper half follows the
% unreviewed draft manuscripts/CJ-0001-upper-bound/main.tex with corrections.
% Plain LaTeX with standard packages only (amsart, amsmath, amssymb, amsthm,
% hyperref).
% Revision 2 (session 14, same day): defects of the two AI reviews of 8df6bea
% (reviews/2026-11-27-claude-{opus,sonnet}-5-5-manuscript.md) corrected;
% Lemma "the star" (lem:starshape) added, so that the bounds no longer use the
% log-concavity theorem of [AAB] or Collins's Theorem 1; Lemma R1 removed.
% Revision 3 (session 15, same day): the findings of the independent review of
% revision 2 by an OpenAI model through Codex
% (reviews/2026-11-27-codex-gpt-astra-6-manuscript-r2.*) that this session
% judged justified, all six: the spectrum of the subdivision corrected and the
% derivation of Theorem T10 from [AMW] made a proved remark; the golden ratio
% defined; the Graham--Lovasz peak written with the floor; the open case of
% peak-maximality of the path stated exactly; R(K,nu) defined; the decimal
% estimate at the end of Lemma S44 replaced by an exact one.
\documentclass[11pt]{amsart}
\usepackage{amsmath,amssymb,amsthm}
\PassOptionsToPackage{hyphens}{url}
\usepackage[colorlinks=true,allcolors=blue]{hyperref}

\newtheorem{theorem}{Theorem}[section]
\newtheorem{lemma}[theorem]{Lemma}
\newtheorem{proposition}[theorem]{Proposition}
\newtheorem{corollary}[theorem]{Corollary}
\newtheorem*{maintheorem}{Main Theorem}
\theoremstyle{definition}
\newtheorem{definition}[theorem]{Definition}
\newtheorem{conjecture}[theorem]{Conjecture}
\newtheorem{question}[theorem]{Question}
\theoremstyle{remark}
\newtheorem{remark}[theorem]{Remark}

\newcommand{\tr}{\operatorname{tr}}
\newcommand{\Tr}{\operatorname{Tr}}
\newcommand{\adj}{\operatorname{adj}}
\newcommand{\diag}{\operatorname{diag}}
\newcommand{\frc}[1]{\{#1\}}
\newcommand{\K}{\mathcal K}
\newcommand{\one}{\mathbf 1}
\newcommand{\T}{\mathsf T}

\title[The peak of the distance characteristic polynomial of a tree]{On the
peak of the coefficients of the distance characteristic polynomial of a
tree: a candidate proof of the Collins--Shor conjecture\\[0.5ex]
\normalfont\small (draft prepared with AI assistance, for independent review)}
\author[Conjectura project (draft)]{Conjectura project (authorship not yet
determined; see the attribution section)}
\address{Conjectura project, directed by Gustavo Bernardo Brito, Universidade
Federal de Goi\'as (UFG), Brazil}
\date{Draft of 27 November 2026}

\emergencystretch=3em
\begin{document}

\begin{abstract}
Let $T$ be a tree with $n\ge3$ vertices and distance matrix $D$, and let
$d_0,\dots,d_{n-2}$ be the normalized coefficients
$d_k=2^{k-(n-2)}|\delta_k|$ of $\det(D-xI)=\sum_k\delta_kx^k$. Graham and
Lov\'asz conjectured that this sequence is unimodal with peak at
$\lfloor n/2\rfloor$;
Collins showed that the peak of the path is near $n-n/\sqrt5$, and Aalipour
et al.\ proved that the sequence is
log-concave, hence unimodal, and recorded the conjecture of Collins and Shor
that the peak lies between $\lfloor n/2\rfloor$ and $\lceil n-n/\sqrt5\rceil$.
We give complete written proofs that every index $p$ with $d_p=\max_kd_k$
satisfies
\[
\lfloor n/2\rfloor\ \le\ p\ <\ n-n/\sqrt5+3/5 ,
\]
so in particular $p\le\lceil n-n/\sqrt5\rceil$. The lower bound comes from a
coefficientwise comparison with the star: $\ell_T-\ell_{S_n}$ lies in the cone
generated by the polynomials $x^e(1+x)^i$ with $2e+i\ge n-1$, where $\ell_T$
is the normalized distance polynomial; the proof is an induction on rooted
trees in which the only negative term becomes a sum of squares of differences
of distances. The upper bound comes from the spectrum: $\ell_T=4(\alpha-x)q(x)$
with $q$ real-rooted, the peak of $(d_k)$ is at most the mode of $q$, the mode
is controlled by the mean (Darroch's rule), the mean is at most
$n-\tr(I+L)^{-1}$ by interlacing with the Laplacian spectrum, and the path
minimizes $\tr(I+L)^{-1}$ among trees. Both bounds are attained: the lower
by the star for every $n$, the upper by the path for infinitely many $n$.
This text is a candidate proof that has been checked by AI reviews and exact
computation only; it has not been read by a mathematician.
\end{abstract}

\maketitle

\begin{center}
\fbox{\parbox{0.93\textwidth}{\small
\textbf{Draft prepared with AI assistance, for independent review. Not for
distribution.} The mathematics of this draft was produced by AI systems
(Claude, models \nolinkurl{claude-fable-5-2} and \nolinkurl{claude-opus-5-5},
Anthropic) in the Conjectura research project, November 2026. Its parts were
reviewed only by AI reviewers run in fresh contexts (models of the same
vendor); \textbf{AI review is not peer review}, and \textbf{no mathematician
has read the proofs}. The first version of this document was reviewed as a
whole by two AI reviewers of the same vendor; revision~2 was rechecked by
one of them and reviewed as a whole by an AI model of another vendor
(OpenAI, through Codex), all without finding a gap; this revision~3
corrects the defects that the last review reported, and has been rechecked
only as stated in Section~\ref{sec:status}.
The status of every statement is given in Section~\ref{sec:status}; novelty
is not claimed beyond what is stated in Section~\ref{sec:literature}. Nothing
has been submitted, posted or sent; every outside use is decided by the
project director.}}
\end{center}

\tableofcontents

\section{Introduction}\label{sec:intro}

Let $T$ be a tree with vertex set $V$, $|V|=n$, and let $D$ be its distance
matrix. Graham and Lov\'asz~\cite{GL78} showed that the coefficients
$\delta_k$ of $\det(D-xI)=\sum_k\delta_kx^k$ satisfy
$(-1)^{n-1}\delta_k>0$ for $0\le k\le n-2$ (in the form
\cite[Fact~1.1]{AAB}), so that the \emph{normalized coefficients}
\[
d_k=d_k(T)=2^{k-(n-2)}\,|\delta_k|\qquad(0\le k\le n-2)
\]
are positive. Graham and Lov\'asz~\cite[p.~83]{GL78} observed that these
quantities appear to be unimodal ``with the maximum value occurring for
$k=[n/2]$'' (the integer part, $\lfloor n/2\rfloor$); Collins recalls this
conjecture with peak at $k=\tfrac12n$ \cite[p.~29]{Col89}. Collins~\cite{Col89} computed the sequence for the star and for
the path, showed that the star peaks at $\lfloor n/2\rfloor$ but the path
near $n-n/\sqrt5$ (\cite[Theorems~1 and~3]{Col89}), and stated, as
Conjecture~9 of~\cite{Col89}, attributed there to Peter Shor: ``The
coefficients of $\mathrm{CPD}(T)$ for any tree $T$ with $n$ vertices are
unimodal with peak between $\tfrac12n$ and $n(1-1/\sqrt5)$.'' Aalipour,
Abiad, Berikkyzy, Hogben, Kenter, Lin and Tait~\cite{AAB} proved that the
sequence is log-concave, hence unimodal (\cite[Theorem~2.1]{AAB}), proved that
its peak is at most $\lceil 2n/3\rceil$ (\cite[Theorem~3.3]{AAB}), and
restated the location of the peak as follows.

\begin{conjecture}[{\cite[Conjecture~1.2]{AAB}}; Collins--Shor]\label{conj:CS}
``The [normalized] coefficients of the distance characteristic polynomial for
any tree $T$ with $n$ vertices are unimodal with peak between
$\lfloor n/2\rfloor$ and $\lceil n-n/\sqrt5\rceil$.''
\end{conjecture}

(The bracketed word is in the source.) For odd $n$ the literal bound
$\tfrac12n$ of~\cite{Col89} would contradict~\cite[Theorem~1]{Col89}, which
places the peak of the star at $\tfrac12(n-1)$; the floor in
Conjecture~\ref{conj:CS} is the consistent reading, and we prove that one.
Neither source says how a tie between two maximal coefficients is to be read.
We prove the strong reading, in which \emph{every} index $p$ with
$d_p=\max_kd_k$ (a \emph{maximizing index}) is bounded. As in~\cite{AAB},
whose standing convention is a tree of order at least three, we take
$n\ge3$; for $n=2$ the sequence is the single term $d_0$, and the lower
bound $\lfloor n/2\rfloor=1$ would fail.

\begin{maintheorem}
Let $T$ be a tree with $n\ge3$ vertices. Every maximizing index $p$ of
$d_0(T),\dots,d_{n-2}(T)$ satisfies
\[
\Bigl\lfloor\frac n2\Bigr\rfloor\ \le\ p\ \le\ \Bigl\lceil n-\frac n{\sqrt5}-\frac25\Bigr\rceil\ \le\ \Bigl\lceil n-\frac n{\sqrt5}\Bigr\rceil .
\]
Both outer bounds are attained: the lower one by the star $S_n$ for every
$n\ge3$ (its unique maximizing index is $\lfloor n/2\rfloor$), the upper one
by the path $P_n$ for infinitely many $n$ (this uses Collins's formula for
the path, (C) of Section~\ref{sec:tools}).
\end{maintheorem}

Together with the unimodality of $(d_k)$, \cite[Theorem~2.1]{AAB}, this is
Conjecture~\ref{conj:CS}. The proof of the two bounds does not use that
theorem. The
lower bound is Theorem~\ref{thm:lower}(c), the upper bound
Theorem~\ref{thm:upper}; attainment is Remark~\ref{rem:star-peak} and
Corollary~\ref{cor:pathceil}.

\subsection*{The two halves}
The two halves are proved by unrelated methods and can be read independently
(Sections~\ref{sec:cones}--\ref{sec:lowerproof} and
Sections~\ref{sec:spectral}--\ref{sec:path}); they share only
Section~\ref{sec:prelim}.

\emph{Lower half.} Put $\ell_T(x)=-2^{-(n-2)}\det(2xI-D)$, whose coefficient
of $x^k$ is $d_k$ for $k\le n-2$. For an integer $d$ let $\K_d$ be the convex
cone of polynomials spanned by $x^e(1+x)^i$ with $e,i\ge0$ and $2e+i\ge d$.
A short binomial computation (Lemma~\ref{lem:A}) shows that every element $f$
of $\K_{2\mu-1}$ satisfies $[x^\mu]f\ge[x^{\mu-1}]f$. We prove
(Theorem~\ref{thm:lower}(b)) that
\[
\ell_T-\ell_{S_n}\in\K_{n-1}\qquad\text{for every tree $T$ of order $n$,}
\]
which gives $d_k(T)\ge d_k(S_n)$ for every $k$ and, with $\mu=\lfloor n/2\rfloor$,
$d_\mu(T)-d_{\mu-1}(T)\ge d_\mu(S_n)-d_{\mu-1}(S_n)>0$. The proof starts from
the Graham--Lov\'asz formula $\ell_T=(n-1)\det(I+xL)-x\tau^\T\adj(I+xL)\tau$
(Lemma~\ref{lem:L1}), builds a tree one leaf at a time
(Corollary~\ref{cor:leaf}), and reduces everything to one statement about
rooted trees, $R(A,e_u)\in\K_{|A|+1}$ (Theorem~\ref{thm:T5}), proved by
induction on rooted trees. The induction closes because the one negative
term of the recursion, $-2G_1G_2$, combines with two other terms into
$x^4\sum(\delta_v-\delta_{v'})^2a_va_{v'}$, a sum of squares of differences of
distances to the roots (Proposition~\ref{prop:S35}).

\emph{Upper half.} By the Graham--Lov\'asz inverse and Weyl's interlacing,
$\ell_T=4(\alpha-x)q(x)$ where $q$ has only negative zeros
(Lemma~\ref{lem:S40}). The peak of $(d_k)$ is at most the mode of $q$
(Lemma~\ref{lem:S41}); Darroch's rule bounds the mode of $q$ by its mean plus
one; interlacing bounds the mean by $n-\tr(I+L)^{-1}$ (Lemma~\ref{lem:S42});
and $\tr(I+L)^{-1}$ is minimized over trees of order $n$ by the path
(Theorem~\ref{thm:T10}), where it equals
$n/\sqrt5+2/5+O(n\varphi^{-4n})$, with $\varphi=(1+\sqrt5)/2$ the golden
ratio (Lemma~\ref{lem:S46}).

\subsection*{What is new and what is known}
Section~\ref{sec:literature} gives the details and the coverage of the
searches made. In short: the formula of Lemma~\ref{lem:L1}, the path formula
for the off-diagonal cofactors (Lemma~\ref{lem:path}), the polynomials of the
star and of the path, Darroch's rule and the interlacing of
Lemmas~\ref{lem:S40} and~\ref{lem:S42} are known; Theorem~\ref{thm:T10} is
a consequence of a theorem of Andriantiana, Razanajatovo Misanantenaina and
Wagner~\cite{AMW} on the average size of matchings
(Remark~\ref{rem:AMW}; our proof is independent of theirs); no prior
proof of either half of Conjecture~\ref{conj:CS}, and no prior statement of
the coefficientwise minimality of the star (Theorem~\ref{thm:lower}(b)), was
found within the coverage of the searches. That is not evidence of
originality.

\subsection*{Conventions}
Polynomials are in $\mathbb R[x]$; $[x^k]f$ is the coefficient of $x^k$ in
$f$; $\binom ik=0$ unless $0\le k\le i$. For a matrix $N$ indexed by $V$ and
$S\subseteq V$, $N[S]$ is the principal submatrix on $S$, with
$\det N[\emptyset]=1$; $\adj N$ is the adjugate, $\adj(N)N=N\adj(N)=(\det N)I$;
$E_{uu}$ is the matrix unit and $e_u$ the unit vector at $u$; $\one$ is the
all-ones vector; $\varphi=(1+\sqrt5)/2$. An identity between polynomials in $x$ that holds for all
real $x>0$ holds identically; we use this to prove polynomial identities by
matrix inversion at positive $x$, where the matrices involved are positive
definite.

\section{Preliminaries}\label{sec:prelim}

\subsection{The normalized polynomial}
Let $T$ be a tree with $n\ge1$ vertices, distance matrix $D$, Laplacian
$L=L_T$, degrees $k_v$, and $\tau\in\mathbb Z^V$, $\tau_v=2-k_v$; so
$\one^\T\tau=2n-2(n-1)=2$. Put
\[
\ell_T(x)=-2^{-(n-2)}\det(2xI-D),\qquad M=I+xL,\qquad \Phi_T=\det M .
\]
For $n=1$, $\ell_T=-4x$. The coefficient of $x^n$ in $\ell_T$ is $-4$ and that
of $x^{n-1}$ is $0$ (since $\tr D=0$).

\begin{lemma}[{\cite[Observation~3.1]{AAB}}]\label{lem:ident}
For $n\ge2$ and $0\le k\le n-2$, $[x^k]\ell_T=d_k(T)$.
\end{lemma}

\begin{proof}
$\det(2xI-D)=(-1)^n\det(D-2xI)=(-1)^n\sum_k\delta_k2^kx^k$, so
$[x^k]\ell_T=(-1)^{n-1}2^{k-(n-2)}\delta_k$, and $(-1)^{n-1}\delta_k>0$ for
$0\le k\le n-2$ by \cite{GL78} in the form \cite[Fact~1.1]{AAB}.
\end{proof}

\subsection{Classical results used}\label{sec:tools}
The following are used without proof. For each we say where it is used and
how it was checked in the project.

\begin{enumerate}
\item[(GL)] \emph{Graham--Lov\'asz inverse} \cite[Lemma~1, p.~66]{GL78}: for
a tree with $n\ge2$ vertices, $D^{-1}=-\frac12L+\frac{1}{2(n-1)}\tau\tau^\T$.
Used in Lemmas~\ref{lem:L1} and~\ref{lem:S40}. Read in a scan of~\cite{GL78}
(two signs illegible in the scan were fixed from other formulas of the same
paper); verified in exact arithmetic for every tree of order $3$ to $17$.
\item[(GP)] \emph{Graham--Pollak determinant} \cite[proof of Theorem~3,
pp.~2509--2511]{GP71}: $\det D=(-1)^{n-1}(n-1)2^{n-2}$ for $n\ge2$. Used in
Lemmas~\ref{lem:L1} and~\ref{lem:S40}. Read at the source by the project's
literature auditor (a scan); also quoted in~\cite{AAB}; consistent with all
exact computations.
\item[(LC)] \emph{Log-concavity} \cite[Theorem~2.1]{AAB}: the sequence
$d_0,\dots,d_{n-2}$ is positive and log-concave. Not used in the proofs below;
it supplies the word ``unimodal'' of Conjecture~\ref{conj:CS}. Read in the
source.
\item[(W)] \emph{Weyl interlacing for a rank-one positive semidefinite
update}: if $B=A+ww^\T$ with $A$ real symmetric of size $n$, then, with
eigenvalues in decreasing order, $\lambda_i(A)\le\lambda_i(B)\le\lambda_{i-1}(A)$
for $2\le i\le n$, and $\lambda_1(A)\le\lambda_1(B)$ \cite{HJ}.
\item[(N)] \emph{Newton's inequalities}: if $\sum_{k=0}^Nq_kx^k$ has positive
coefficients and only real zeros, then
$q_k^2\ge q_{k-1}q_{k+1}(1+1/k)(1+1/(N-k))>q_{k-1}q_{k+1}$ for $1\le k\le N-1$
\cite{HLP}.
\item[(D)] \emph{Darroch's rule} \cite[Theorem~4]{Dar64}, in the form stated
by Pitman~\cite[p.~284]{Pit97}: if a polynomial $P$ with non-negative
coefficients has only real zeros, every index at which its coefficient is
maximal lies at distance less than $1$ from $\mu=P'(1)/P(1)$. Darroch's paper
was not read in the project (its numbering, Theorem~4, is the one given by
the secondary sources); the statement was read in~\cite{Pit97},
including the case of two maximizing indices \cite[eq.~(13)]{Pit97}, and
tested on $20{,}000$ random real-rooted polynomials.
\item[(M)] \emph{M-matrix facts} \cite{BP}: for a connected graph $H$,
$(I+L_H)^{-1}$ exists and is entrywise positive; $(I+L_H)\one=\one$.
\item[(PS)] \emph{Spectrum of the path}: the Laplacian eigenvalues of $P_n$
are $2-2\cos(j\pi/n)$, $0\le j\le n-1$ \cite{BH}.
\item[(C)] \emph{Collins's formula for the path} \cite[Theorem~3 and its
proof, p.~31]{Col89}: for $n\ge3$ and $2\le k\le n$,
$d_{n-k}(P_n)=\binom{n+k-1}{2k-1}\frac{(k-1)n}{k}$. Used only in
Section~\ref{sec:path}. Read in the source; checked against the distance
matrix in exact arithmetic for $3\le n\le30$.
\item[(Kr)] \emph{Kronecker}: for irrational $\zeta$, the fractional parts of
$n\zeta$, $n\ge1$, are dense in $[0,1)$. Used only in
Corollary~\ref{cor:pathceil}.
\end{enumerate}
The Schur complement formulas, the matrix determinant lemma, the
Sherman--Morrison formula and the Poisson kernel are used as elementary
identities. (W), (N), (M), (PS), (Kr) are textbook facts and were not checked
against a particular source.

\subsection{The Laplacian form}

\begin{lemma}[{Graham--Lov\'asz~\cite{GL78}; the Laplacian form}]\label{lem:L1}
For every tree with $n\ge2$ vertices,
$\ell_T=(n-1)\Phi_T-x\,\tau^\T\adj(M)\,\tau$.
\end{lemma}

\begin{proof}
By (GL), $2xD^{-1}-I=-(I+xL)+\frac{x}{n-1}\tau\tau^\T$. By the matrix
determinant lemma $\det(A+uv^\T)=\det A+v^\T\adj(A)u$ with $A=-M$
($\det A=(-1)^n\Phi_T$, $\adj A=(-1)^{n-1}\adj M$),
$\det(2xD^{-1}-I)=(-1)^n\bigl[\Phi_T-\frac{x}{n-1}\tau^\T\adj(M)\tau\bigr]$.
Multiply by $\det D=(-1)^{n-1}(n-1)2^{n-2}$ (GP) to get
$\det(2xI-D)=-2^{n-2}\bigl[(n-1)\Phi_T-x\tau^\T\adj(M)\tau\bigr]$.
\end{proof}

The formula also holds for $n=2$ by direct computation:
$\ell=1-4x^2$, $\Phi=1+2x$, $\adj M=\left(\begin{smallmatrix}1+x&x\\x&1+x\end{smallmatrix}\right)$,
$\tau=(1,1)$, and $(1+2x)-x(2+4x)=1-4x^2$; Lemma~\ref{lem:L1} covers this
case too.

\part*{The lower bound}

\section{The cones \texorpdfstring{$\K_d$}{K\_d}}\label{sec:cones}

\begin{definition}
For an integer $d$, $\K_d$ is the set of finite linear combinations with
non-negative real coefficients of the polynomials $x^e(1+x)^i$, $e,i\ge0$
integers, $2e+i\ge d$.
\end{definition}

Directly from the definition:
\begin{enumerate}
\item[(K1)] elements of $\K_d$ have non-negative coefficients; $\K_d$ is closed
under addition; $\K_d=\K_0$ for $d\le0$, and $\K_0$ is the set of all
polynomials with non-negative coefficients;
\item[(K2)] $\K_{d+1}\subseteq\K_d$, $\K_d\K_{d'}\subseteq\K_{d+d'}$,
$(1+x)\K_d\subseteq\K_{d+1}$, $x\K_d\subseteq\K_{d+2}$; hence
$(1+4x)\K_d=((1+x)+3x)\K_d\subseteq\K_{d+1}$ and
$x(2+5x)\K_d=(2x(1+x)+3x^2)\K_d\subseteq\K_{d+3}$.
\end{enumerate}
For $\mu\ge1$ and a polynomial $f$ put $\lambda_\mu(f)=[x^\mu]f-[x^{\mu-1}]f$.

\begin{lemma}\label{lem:A}
If $\mu\ge1$ and $f\in\K_{2\mu-1}$, then $\lambda_\mu(f)\ge0$.
\end{lemma}

\begin{proof}
$\lambda_\mu$ is linear, so it suffices to take $f=x^e(1+x)^i$ with
$2e+i\ge2\mu-1$. Let $a=\mu-e$. Then $[x^\mu]f=\binom ia$ and
$[x^{\mu-1}]f=\binom i{a-1}$. If $a\le0$, $\lambda_\mu(f)=\binom ia\ge0$. If
$a\ge1$, then $i\ge2\mu-1-2e=2a-1\ge a$, so
$\binom ia=\binom i{a-1}\frac{i-a+1}{a}$ and
$\lambda_\mu(f)=\binom i{a-1}\frac{i-2a+1}{a}\ge0$.
\end{proof}

\section{Minors of \texorpdfstring{$I+x(L+\diag\nu)$}{I+x(L+diag nu)}}\label{sec:minors}

A \emph{weighted tree} $(K,\nu)$ is a tree $K$ with vertex set $U$, $|U|=j\ge1$,
and $\nu\in\mathbb Z_{\ge0}^U$. Put
\[
M(K,\nu)=I+x(L_K+\diag\nu),\qquad \Phi(K,\nu)=\det M(K,\nu).
\]
For real $x>0$, $M(K,\nu)$ and all its principal submatrices are positive
definite ($L_K$ is positive semidefinite). If $j=1$, $\Phi(K,\nu)=1+\nu x$.

\begin{lemma}\label{lem:blocks}
Let $(K,\nu)$ be a weighted tree and $Q\subseteq U$. Then $M(K,\nu)[Q]$ is
block diagonal over the components $K'$ of the forest $K[Q]$, with blocks
$M(K',\nu|_{K'}+\beta_{K'})$, where $\beta_{K'}(w)$ is the number of
neighbours of $w$ outside $Q$. If $Q=U\setminus V(P)$ for a non-empty path
$P$ of $K$, then every $\beta_{K'}$ is a unit vector $e_{a_{K'}}$.
\end{lemma}

\begin{proof}
The diagonal entry of $M(K,\nu)$ at $w$ is $1+x(k_w+\nu_w)$, and $k_w$ is the
degree of $w$ in $K'$ plus $\beta_{K'}(w)$; the off-diagonal entries are $-x$
on the edges and $0$ elsewhere, and no edge joins two components of $K[Q]$.
If $Q=U\setminus V(P)$, each component $K'$ is joined to $V(P)$ by exactly one
edge (two would close a cycle with a segment of $P$), so $\beta_{K'}=e_{a_{K'}}$,
$a_{K'}$ the end of that edge in $K'$.
\end{proof}

\begin{lemma}\label{lem:K}
Let $(K,\nu)$ be a weighted tree on $j\ge1$ vertices with $\nu\ne0$. Then
$\Phi(K,\nu)\in\K_j$. Moreover, if $w$ is a leaf of $K$ with neighbour $u$ and
$\nu'=\nu|_{U\setminus w}$, then
\begin{equation}\label{eq:K3}
\Phi(K,\nu)=(1+x)\Phi(K-w,\nu')+x\det M(K,\nu)[U\setminus\{w,u\}]+\nu_wx\,\Phi(K-w,\nu'+e_u).
\end{equation}
\end{lemma}

\begin{proof}
Let $N=M(K,\nu)$. Expanding $\det N$ along the row of $w$ (diagonal entry
$1+(1+\nu_w)x$, entry $-x$ at $u$):
$\det N=(1+(1+\nu_w)x)\det N[U\setminus w]-x^2\det N[U\setminus\{w,u\}]$.
Now $N[U\setminus w]=M(K-w,\nu')+xE_{uu}$ (Lemma~\ref{lem:blocks}), and the
determinant is affine in the $(u,u)$ entry, so
$\det N[U\setminus w]=\Phi(K-w,\nu')+x\det N[U\setminus\{w,u\}]=\Phi(K-w,\nu'+e_u)$.
Substituting gives~\eqref{eq:K3}.

The membership is proved by induction on $j$. For $j=1$,
$\Phi=1+\nu x=(1+x)+(\nu-1)x$ with $\nu\ge1$, and $1+x,x\in\K_1$. Let $j\ge2$.
A tree with at least two vertices has two leaves; choose a leaf $w$ such that
$\nu'\ne0$ (if $\nu$ is supported on one vertex, take a leaf other than that
vertex). In~\eqref{eq:K3}: $\Phi(K-w,\nu')\in\K_{j-1}$ by induction, so the
first term is in $\K_j$. By Lemma~\ref{lem:blocks} (with $P=(w,u)$),
$\det N[U\setminus\{w,u\}]$ is a product of determinants $\Phi(K',\cdot)$ with
non-zero weights over components $K'$ with $j-2$ vertices in all; it lies in
$\K_{j-2}$ by induction and (K2) (it is $1\in\K_0$ if $j=2$); times $x$ it is in
$\K_j$. Finally $\Phi(K-w,\nu'+e_u)\in\K_{j-1}$ by induction and
$x\K_{j-1}\subseteq\K_{j+1}\subseteq\K_j$.
\end{proof}

\begin{lemma}\label{lem:path}
Let $(K,\nu)$ be a weighted tree on $j$ vertices, $M=M(K,\nu)$, $v,v'\in U$,
$P$ the path from $v$ to $v'$, $d=\operatorname{dist}(v,v')$ and
$Q=U\setminus V(P)$. Then
\[
\adj(M)_{vv'}=x^d\det M[Q],\qquad \det M[Q]\in\K_{j-d-1}.
\]
\end{lemma}

\begin{proof}
For $d=0$ this is the definition of $\adj(M)_{vv}$. Let $d\ge1$ and $x>0$.
Write $P=(p_0=v,\dots,p_d=v')$. By Lemma~\ref{lem:blocks}, every component of
$K[Q]$ has exactly one vertex adjacent to $P$, and that vertex has exactly one
neighbour on $P$; hence in $M[V(P),Q]\,M[Q]^{-1}M[Q,V(P)]$ the entry
$(p_i,p_{i'})$ vanishes for $i\ne i'$ (it is a sum of terms
$x^2(M[Q]^{-1})_{aa'}$ over $a\sim p_i$, $a'\sim p_{i'}$ with $a,a'$ in the same
component, and then $a=a'$ and $p_i=p_{i'}$). So the Schur complement
$S=M[V(P)]-M[V(P),Q]M[Q]^{-1}M[Q,V(P)]$ differs from $M[V(P)]$ only on the
diagonal, and it is tridiagonal with $-x$ next to the diagonal. Now
$\det M=\det M[Q]\det S$ and $(M^{-1})[V(P)]=S^{-1}$. The $(0,d)$ entry of
$S^{-1}$ is $(-1)^d$ times the determinant of $S$ without row $d$ and
column $0$, divided by $\det S$; that matrix is lower triangular with $-x$ on
its diagonal, so $(S^{-1})_{0d}=x^d/\det S$ and
$\adj(M)_{vv'}=\det M\cdot x^d/\det S=x^d\det M[Q]$.

By Lemma~\ref{lem:blocks}, $\det M[Q]$ is the product of $\Phi(K',\nu|_{K'}+e_{a_{K'}})$
over the components $K'$ of $K[Q]$, each in $\K_{|K'|}$ by Lemma~\ref{lem:K};
by (K2) the product is in $\K_{j-d-1}$ (it is $1$ if $Q=\emptyset$).
\end{proof}

\begin{corollary}\label{cor:phi1}
For every tree $B$ on $b\ge1$ vertices and every vertex $c$,
$\Phi(B,e_c)-(1+x)^b\in\K_b$.
\end{corollary}

\begin{proof}
Induction on $b$; for $b=1$ the difference is $0$. For $b\ge2$ take a leaf
$w\ne c$ with neighbour $y$. By~\eqref{eq:K3} with $\nu=e_c$, $\nu_w=0$:
$\Phi(B,e_c)-(1+x)^b=(1+x)\bigl[\Phi(B-w,e_c)-(1+x)^{b-1}\bigr]+x\det M(B,e_c)[B\setminus\{w,y\}]$.
The first term is in $\K_b$ by induction and (K2); the minor is the one of
Lemma~\ref{lem:path} for the path $(w,y)$, so it lies in $\K_{b-2}$, and
$x\K_{b-2}\subseteq\K_b$.
\end{proof}

\section{Rooted trees}\label{sec:rooted}

\subsection{Definitions}
A \emph{rooted tree} $(A,u)$ is a tree $A$ with $a=|A|\ge1$ vertices and a
vertex $u$, the root. Degrees $k_v$ and $\tau_v=2-k_v$ are taken in $A$. For
an integer $w\ge0$ (the \emph{root weight}) put
\[
M_w=I+x(L_A+wE_{uu}),\quad \Phi_w=\det M_w,\quad p^{(w)}=\one+x\tau-wx\,e_u,\quad
\pi_w=e_u^\T\adj(M_w)\,p^{(w)},
\]
\[
R_w=R_w(A,u)=(w+a-1)(1+x)\Phi_w-p^{(w)\T}\adj(M_w)\,p^{(w)} .
\]
So $M_w=M(A,we_u)$ in the notation of Section~\ref{sec:minors}, and
$p^{(w)}_v=1+(2-k_v-w[v=u])x$. The polynomial $R_0$ does not depend on the
root.

The \emph{branches} of $(A,u)$ are the rooted trees $(A_i,c_i)$,
$1\le i\le k_u$, where $A_1,\dots,A_{k_u}$ are the components of $A-u$ and
$c_i\in A_i$ is the neighbour of $u$. Put
\[
\Psi=\prod_i\Phi_1(A_i,c_i),\qquad
\Omega=\sum_iR_1(A_i,c_i)\prod_{l\ne i}\Phi_1(A_l,c_l)
\]
($\Psi=1$ and $\Omega=0$ if $a=1$), and
\[
H=(\Phi_0-\Psi)/x,\qquad G=(1+2x)\Psi-\pi_0 .
\]
By Lemma~\ref{lem:blocks} with $Q=A\setminus u$, $M_w[A\setminus u]$ is the
direct sum of the matrices $M_1(A_i,c_i)$; so $\Psi=\det M_w[A\setminus u]=\adj(M_w)_{uu}$
for every $w$. $H$ is a polynomial because $\Phi_0$ and $\Psi$ have constant
term $1$. The polynomials $(\Psi,H,G,\Omega,R_0)$ are the \emph{state} of
$(A,u)$.

Three operations build all rooted trees from single vertices:
\begin{itemize}
\item \emph{plant}: from $(A,u)$, form $(\hat A,r)$ by adding a new vertex $r$
adjacent to $u$ only, with root $r$;
\item \emph{merge}: from $(A_1,u_1)$ and $(A_2,u_2)$ on disjoint vertex sets,
identify $u_1$ with $u_2$; the identified vertex is the root;
\item every rooted tree $(A,u)$ with $a\ge2$ is the merge of the planted
trees $(\hat A_i,r_i)$ of its branches (the vertices $r_i$ all identified
with $u$).
\end{itemize}

\subsection{The basic identity}

\begin{lemma}\label{lem:star}
For every rooted tree $(A,u)$ and every $w\ge0$,
\[
\Psi R_w=\Phi_w\bigl[w(1+x)\Psi+\Omega\bigr]-\pi_w^2 .
\]
Moreover, writing $M_w=\left(\begin{smallmatrix}\Delta&b^\T\\b&Y\end{smallmatrix}\right)$
with $u$ first ($\Delta=1+x(k_u+w)$, $b=-x\sum_ie_{c_i}$, $Y=M_w[A\setminus u]$)
and $q$ for the restriction of $p^{(w)}$ to $A\setminus u$,
\begin{equation}\label{eq:blockform}
\Phi_w=\Psi\Delta-b^\T\adj(Y)b,\qquad \pi_w=\Psi p^{(w)}_u-b^\T\adj(Y)q .
\end{equation}
\end{lemma}

\begin{proof}
If $a=1$: $\Psi=1$, $\Omega=0$, $\Phi_w=1+wx$, $\adj(M_w)=(1)$,
$p^{(w)}=\pi_w=1+(2-w)x$, and $R_w=w(1+x)(1+wx)-\pi_w^2$, which is the claim.
Let $a\ge2$ and $x>0$. Let $s=\Delta-b^\T Y^{-1}b$. By the Schur complement
formulas $\Phi_w=\Psi s$, $(M_w^{-1})_{uu}=1/s$, the rest of row $u$ of
$M_w^{-1}$ is $-b^\T Y^{-1}/s$, and
\[
p^\T M_w^{-1}p=q^\T Y^{-1}q+(p_u-b^\T Y^{-1}q)^2/s\qquad(p=p^{(w)}).
\]
Hence $\pi_w=\Phi_we_u^\T M_w^{-1}p=\Psi(p_u-b^\T Y^{-1}q)$, which
with $\Phi_w=\Psi s$ is~\eqref{eq:blockform}, and multiplying the last display by
$\Phi_w\Psi$:
$\Psi\,p^\T\adj(M_w)p=\Phi_w\,q^\T\adj(Y)q+\pi_w^2$.

$Y$ is the direct sum of $Y_i=M_1(A_i,c_i)$, so
$\adj(Y)=\bigoplus_i\adj(Y_i)\prod_{l\ne i}\det Y_l$. The restriction $q_i$ of
$q$ to $A_i$ is the vector $p^{(1)}$ of $(A_i,c_i)$: for $v\in A_i\setminus c_i$
the degree is the same in $A$ and $A_i$, and $c_i$ has degree one more in
$A$, which is the term $-x$ of the root weight $1$. By the definition of
$R_1(A_i,c_i)$, $q_i^\T\adj(Y_i)q_i=a_i(1+x)\Phi_1(A_i,c_i)-R_1(A_i,c_i)$, with
$a_i=|A_i|$, $\sum_ia_i=a-1$. Therefore
$q^\T\adj(Y)q=(a-1)(1+x)\Psi-\Omega$, and
\[
\Psi R_w=(w+a-1)(1+x)\Phi_w\Psi-\Phi_w\bigl[(a-1)(1+x)\Psi-\Omega\bigr]-\pi_w^2,
\]
which is the claim.
\end{proof}

\begin{proposition}\label{prop:rules}
\textup{(a) One vertex.} If $a=1$: $\Psi=\Phi_0=1$, $H=G=\Omega=0$,
$R_0=-(1+2x)^2$, $\Phi_1=1+x$, $R_1=0$.

\textup{(b) Raising the root weight.} $\Phi_1=\Phi_0+x\Psi$,
$\pi_1=\pi_0-x\Psi$, and
\[
R_1=R_0+(1+2x)^2\Psi+x(1+x)H+x\Omega-2xG .
\]

\textup{(c) Plant.} The state of $(\hat A,r)$ is
\[
\hat\Psi=\Phi_1,\quad \hat\Phi_0=(1+x)\Phi_0+x\Psi,\quad \hat H=\Phi_0,\quad
\hat G=x(xH+G),\quad \hat\Omega=R_1,\quad \hat R_0=(1+x)R_0+x\Omega .
\]

\textup{(d) Merge.} If $(A,u)$ is the merge of $(A_1,u_1)$ and $(A_2,u_2)$,
with states indexed by $1,2$: $\Psi=\Psi_1\Psi_2$; $X=X_1\Psi_2+X_2\Psi_1$ for
$X\in\{H,G,\Omega\}$; and
\[
R_0=R_{0,1}\Psi_2+R_{0,2}\Psi_1+x(H_1\Omega_2+H_2\Omega_1)+(1+2x)^2\Psi_1\Psi_2-2G_1G_2 .
\]
\end{proposition}

\begin{proof}
(a) Direct: $p^{(0)}=1+2x$, $p^{(1)}=1+x$, $R_1=(1+x)^2-(1+x)^2=0$.

(b) $M_1=M_0+xE_{uu}$; the determinant is affine in the $(u,u)$ entry, with
coefficient the cofactor $\Psi$; so $\Phi_1=\Phi_0+x\Psi$. The cofactors of
the entries of row $u$ do not involve row $u$, so row $u$ of the adjugate is
the same for $M_0$ and $M_1$; with $p^{(1)}=p^{(0)}-xe_u$ this gives
$\pi_1=\pi_0-x\Psi$. $\Psi$ and $\Omega$ do not depend on $w$. Subtracting the
identities of Lemma~\ref{lem:star} for $w=1$ and $w=0$:
\[
\Psi(R_1-R_0)=(\Phi_0+x\Psi)(1+x)\Psi+x\Psi\Omega-\bigl[(\pi_0-x\Psi)^2-\pi_0^2\bigr],
\]
so, dividing by the non-zero polynomial $\Psi$,
$R_1-R_0=(1+x)\Phi_0+x(1+x)\Psi+x\Omega+2x\pi_0-x^2\Psi$. Substitute
$\Phi_0=\Psi+xH$ and $\pi_0=(1+2x)\Psi-G$; the coefficient of $\Psi$ is
$(1+x)+x(1+x)+2x(1+2x)-x^2=(1+2x)^2$.

(c) The only branch of $(\hat A,r)$ is $(A,u)$; this gives $\hat\Psi=\Phi_1$
and $\hat\Omega=R_1$. With $r$ first,
$\hat M_0=\left(\begin{smallmatrix}1+x&-xe_u^\T\\-xe_u&M_1\end{smallmatrix}\right)$
($u$ has gained one neighbour). Expanding,
$\hat\Phi_0=(1+x)\Phi_1-x^2\adj(M_1)_{uu}=(1+x)(\Phi_0+x\Psi)-x^2\Psi=(1+x)\Phi_0+x\Psi$,
and $\hat H=(\hat\Phi_0-\hat\Psi)/x=\Phi_0$. For $x>0$ the row $r$ of
$\hat M_0^{-1}$ outside $r$ is $x e_u^\T M_1^{-1}/s$ with
$s=\hat\Phi_0/\Phi_1$, so row $r$ of $\adj(\hat M_0)$ is
$(\Phi_1,\ x\,e_u^\T\adj(M_1))$, and $e_u^\T\adj(M_1)=e_u^\T\adj(M_0)$ as in (b).
The vector $\hat p^{(0)}$ is $1+x$ at $r$ and $p^{(0)}-xe_u=p^{(1)}$ on $A$.
Hence $\hat\pi_0=(1+x)\Phi_1+x\pi_1$ and
\[
\hat G=(1+2x)\Phi_1-\hat\pi_0=x(\Phi_1-\pi_1)=x(\Phi_0-\pi_0+2x\Psi)=x(xH+G).
\]
For $\hat R_0$ use Lemma~\ref{lem:star} for $(\hat A,r)$ with $w=0$:
$\Phi_1\hat R_0=\hat\Phi_0R_1-\hat\pi_0^2$. Write
$\hat\Phi_0=(1+x)\Phi_1-x^2\Psi$, $\hat\pi_0=(1+x)\Phi_1+x\pi_1$, and use
Lemma~\ref{lem:star} for $(A,u)$ with $w=1$,
$\Psi R_1=\Phi_1\mathcal A-\pi_1^2$ with $\mathcal A=(1+x)\Psi+\Omega$:
\[
\hat\Phi_0R_1-\hat\pi_0^2=(1+x)\Phi_1R_1-x^2(\Phi_1\mathcal A-\pi_1^2)-(1+x)^2\Phi_1^2-2x(1+x)\Phi_1\pi_1-x^2\pi_1^2,
\]
so $\hat R_0=(1+x)R_1-x^2\mathcal A-(1+x)^2\Phi_1-2x(1+x)\pi_1$. Insert (b) in
the form $R_1=R_0+(1+x)\Phi_0+x(1+x)\Psi+x\Omega+2x\pi_0-x^2\Psi$, together
with $\Phi_1=\Phi_0+x\Psi$, $\pi_1=\pi_0-x\Psi$. The terms in $\Phi_0$ cancel
($(1+x)^2-(1+x)^2$), so do those in $\pi_0$ ($2x(1+x)-2x(1+x)$); the terms in
$\Omega$ give $x(1+x)-x^2=x$; and the coefficient of $\Psi$ is
$(1+x)\bigl[x(1+x)-x^2\bigr]-x^2(1+x)-x(1+x)^2+2x^2(1+x)=(1+x)\bigl[x-x^2-x-x^2+2x^2\bigr]=0$.
So $\hat R_0=(1+x)R_0+x\Omega$.

(d) The branches of the merged tree are those of the two trees; this gives
$\Psi$ and $\Omega$. In~\eqref{eq:blockform} with $w=0$, the merged tree has
$k_u=k_{u_1}+k_{u_2}$, $Y=Y_1\oplus Y_2$, $b=(b_1,b_2)$ and $q=(q_1,q_2)$
(degrees off the root do not change), so $\adj(Y)=\Psi_2\adj(Y_1)\oplus\Psi_1\adj(Y_2)$,
and each of $\beta=b^\T\adj(Y)b$, $B=b^\T\adj(Y)q$ and $k_u\Psi$ has the form
$X_1\Psi_2+X_2\Psi_1$. Since $xH=\Phi_0-\Psi=xk_u\Psi-\beta$ and
$G=(1+2x-p_u)\Psi+B=xk_u\Psi+B$, so do $H$ and $G$.

For $R_0$, Lemma~\ref{lem:star} with $w=0$ reads $\Psi R_0=\Phi_0\Omega-\pi_0^2$.
For the merged tree
$\Phi_0=\Psi+xH=\Phi_{0,1}\Psi_2+\Phi_{0,2}\Psi_1-\Psi_1\Psi_2$ and
$\pi_0=(1+2x)\Psi-G=\pi_{0,1}\Psi_2+\pi_{0,2}\Psi_1-c\Psi_1\Psi_2$, $c=1+2x$.
Expanding $\Psi_1\Psi_2R_0=\Phi_0\Omega-\pi_0^2$, the terms
$\Psi_2^2(\Phi_{0,1}\Omega_1-\pi_{0,1}^2)=\Psi_2^2\Psi_1R_{0,1}$ and the
symmetric one appear, and every other term is divisible by $\Psi_1\Psi_2$;
after division,
\[
R_0=R_{0,1}\Psi_2+R_{0,2}\Psi_1+(\Phi_{0,1}-\Psi_1)\Omega_2+(\Phi_{0,2}-\Psi_2)\Omega_1
-2\pi_{0,1}\pi_{0,2}+2c(\pi_{0,1}\Psi_2+\pi_{0,2}\Psi_1)-c^2\Psi_1\Psi_2 .
\]
The last three terms equal $-2(\pi_{0,1}-c\Psi_1)(\pi_{0,2}-c\Psi_2)+c^2\Psi_1\Psi_2=-2G_1G_2+c^2\Psi_1\Psi_2$,
and $\Phi_{0,i}-\Psi_i=xH_i$.
\end{proof}

\subsection{The leaf recursion}

For a tree $T$ and a vertex $u$, write $\Psi_u,\Omega_u$ for the polynomials
$\Psi,\Omega$ of $(T,u)$ and put $r_u(T)=(1+4x)\Psi_u+\Omega_u$.

\begin{corollary}\label{cor:leaf}
\textup{(a)} For every tree $T$ with $n\ge1$ vertices,
$R_0(T)=x\ell_T-(1+4x)\Phi_T$.

\textup{(b)} If $T^+$ is $T$ with a new pendant vertex attached at $u$, then
$\ell_{T^+}=(1+x)\ell_T+r_u(T)$.

\textup{(c)} If $T_1\subset T_2\subset\dots\subset T_n=T$, $|T_k|=k$, and
$T_{k+1}$ is $T_k$ with a pendant vertex attached at $u_k$, then
\[
\ell_T=-4x(1+x)^{n-1}+\sum_{k=1}^{n-1}(1+x)^{n-1-k}\,r_{u_k}(T_k).
\]

\textup{(d)} For the star, $\ell_{S_n}=(n-1)(1+4x)(1+x)^{n-2}-4x(1+x)^{n-1}=(1+x)^{n-2}\bigl[(n-1)+4(n-2)x-4x^2\bigr]$ $(n\ge2)$.
\end{corollary}

\begin{proof}
(a) For $n=1$ both sides are $-(1+2x)^2$. For $n\ge2$: $M\one=\one$, so
$\adj(M)\one=\Phi_T\one$, and with $p=\one+x\tau$, $\one^\T\tau=2$:
$p^\T\adj(M)p=n\Phi_T+4x\Phi_T+x^2\tau^\T\adj(M)\tau$. Hence
$R_0=(n-1)(1+x)\Phi_T-p^\T\adj(M)p=x\bigl[(n-1)\Phi_T-x\tau^\T\adj(M)\tau\bigr]-(1+4x)\Phi_T$,
and the bracket is $\ell_T$ by Lemma~\ref{lem:L1}.

(b) $T^+$ rooted at the new vertex is the plant of $(T,u)$. By
Proposition~\ref{prop:rules}(c), $R_0(T^+)=(1+x)R_0(T)+x\Omega_u$ and
$\Phi_{T^+}=(1+x)\Phi_T+x\Psi_u$. By (a) for $T^+$ and $T$,
\[
x\ell_{T^+}=(1+x)\bigl[x\ell_T-(1+4x)\Phi_T\bigr]+x\Omega_u+(1+4x)\bigl[(1+x)\Phi_T+x\Psi_u\bigr]
=x\bigl[(1+x)\ell_T+(1+4x)\Psi_u+\Omega_u\bigr].
\]

(c) Apply (b) $n-1$ times, starting from $\ell_{T_1}=-4x$.

(d) Build the star from its centre $c$: every branch of $(T_k,c)$ is a
single vertex with root weight $1$, for which $\Phi_1=1+x$ and $R_1=0$
(Proposition~\ref{prop:rules}(a)). So $r_c(T_k)=(1+4x)(1+x)^{k-1}$, and (c)
gives the formula.
\end{proof}

\section{Distances from the root}\label{sec:distances}

Let $(A,u)$ be a rooted tree, $\delta_v=\operatorname{dist}(u,v)$, $M=M_0$,
$M^1=M_1$, $\Phi^1=\Phi_1$. For $v\ne u$ let $P_{uv}$ be the path from $u$ to
$v$ and
\[
\Pi_{uv}=\det M[A\setminus V(P_{uv})],\qquad a_v=x^{\delta_v-1}\Pi_{uv},
\]
and put
\[
J=\sum_{v\ne u}\delta_va_v,\qquad W=\sum_{v\ne u}\delta_v^2a_v,\qquad
Q^1=\sum_{\{v,v'\}}(\delta_v-\delta_{v'})^2\adj(M^1)_{vv'},
\]
the last sum over unordered pairs of distinct vertices of $A$.

\begin{proposition}\label{prop:S33}
For every rooted tree $(A,u)$ of order $a$:
\begin{enumerate}
\item[(a)] $a_v\in\K_{a+\delta_v-3}$ for $v\ne u$; so $a_v\in\K_{a-2}$, and
$a_v\in\K_{a-1}$ if $\delta_v\ge2$;
\item[(b)] $H=\sum_{v\ne u}a_v\in\K_{a-2}$;
\item[(c)] $G=x\sum_{v\ne u}(\delta_v-1)a_v=x(J-H)$; $G/x\in\K_{a-1}$;
\item[(d)] $\Psi\in\K_{a-1}$;
\item[(e)] $R_1=\Phi^1\sum_v\delta_v^2-\delta^\T\adj(M^1)\delta=Q^1+x^2W$, and
$Q^1\in\K_a$; $R_1$ has degree at most $a+1$.
\end{enumerate}
\end{proposition}

\begin{proof}
(a) $\Pi_{uv}\in\K_{a-\delta_v-1}$ by Lemma~\ref{lem:path}, and
$x^{\delta_v-1}\K_{a-\delta_v-1}\subseteq\K_{a+\delta_v-3}$ by (K2).

(b) $\adj(M)\one=\Phi_0\one$, so $\Phi_0=\sum_v\adj(M)_{uv}=\Psi+\sum_{v\ne u}x^{\delta_v}\Pi_{uv}$
by Lemma~\ref{lem:path}; that is, $xH=x\sum_{v\ne u}a_v$. Membership by (a).

(c) \emph{Claim: $L\delta=\tau-2e_u$.} For $v\ne u$, one neighbour of $v$ is at
distance $\delta_v-1$ from $u$ and $k_v-1$ are at distance $\delta_v+1$, so
$(L\delta)_v=k_v\delta_v-(\delta_v-1)-(k_v-1)(\delta_v+1)=2-k_v$; and
$(L\delta)_u=-k_u=\tau_u-2$. Hence $p^{(0)}=\one+x\tau=\one+(M-I)\delta+2xe_u$,
and since $e_u^\T\adj(M)\one=\Phi_0$, $e_u^\T\adj(M)M\delta=\Phi_0\delta_u=0$,
\[
\pi_0=\Phi_0-\sum_v\delta_v\adj(M)_{uv}+2x\Psi=\Phi_0+2x\Psi-xJ .
\]
So $G=(1+2x)\Psi-\pi_0=-(\Phi_0-\Psi)+xJ=x(J-H)$. Only $v$ with
$\delta_v\ge2$ contribute, and their $a_v$ lie in $\K_{a-1}$ by (a).

(d) $\Psi=\adj(M)_{uu}\in\K_{a-1}$ by Lemma~\ref{lem:path} with $v=v'=u$.

(e) By the claim in (c), $p^{(1)}=\one+x\tau-xe_u=\one+xL\delta+xe_u$, and
since $\delta_u=0$, $M^1(\one+\delta)=\one+xe_u+\delta+xL\delta$; so
$p^{(1)}=M^1(\one+\delta)-\delta$. With $\adj(M^1)M^1=\Phi^1I$,
\[
p^{(1)\T}\adj(M^1)p^{(1)}=\Phi^1(\one+\delta)^\T M^1(\one+\delta)-2\Phi^1\delta^\T(\one+\delta)+\delta^\T\adj(M^1)\delta .
\]
Each edge joins vertices whose distances to $u$ differ by $1$, so
$(\one+\delta)^\T L(\one+\delta)=\delta^\T L\delta=a-1$, and
$(\one+\delta)^\T M^1(\one+\delta)=\sum_v(1+\delta_v)^2+x(a-1)+x$. Hence
$p^{(1)\T}\adj(M^1)p^{(1)}=\Phi^1\bigl[a(1+x)-\sum_v\delta_v^2\bigr]+\delta^\T\adj(M^1)\delta$,
and $R_1=a(1+x)\Phi^1-p^{(1)\T}\adj(M^1)p^{(1)}=\Phi^1\sum_v\delta_v^2-\delta^\T\adj(M^1)\delta$.
Since $M^1\one=\one+xe_u$, $\sum_{v'}\adj(M^1)_{vv'}+x\adj(M^1)_{vu}=\Phi^1$ for
every $v$; substituting this for $\Phi^1$ in the coefficient of
$\delta_v^2$ and using symmetry,
\[
R_1=x\sum_v\delta_v^2\adj(M^1)_{vu}+\tfrac12\sum_{v,v'}(\delta_v-\delta_{v'})^2\adj(M^1)_{vv'} .
\]
The second sum is $Q^1$. By Lemma~\ref{lem:path} for $(A,e_u)$,
$\adj(M^1)_{vu}=x^{\delta_v}\det M^1[A\setminus V(P_{uv})]$ for $v\ne u$, and
this minor is $\Pi_{uv}$ because $M^1$ and $M$ differ only in the entry
$(u,u)$, which is deleted; so the first sum is $x^2W$. For distinct $v,v'$ at
distance $d\ge1$, Lemma~\ref{lem:path} for $(A,e_u)$ gives
$\adj(M^1)_{vv'}\in x^d\K_{a-d-1}\subseteq\K_{a+d-1}\subseteq\K_a$; so
$Q^1\in\K_a$. Finally $\Phi^1$ has degree at most $a$, the entries of
$\adj(M^1)$ at most $a-1$ and those of $p^{(1)}$ at most $1$.
\end{proof}

\begin{proposition}\label{prop:S34}
For a rooted tree $(A,u)$ put $E=R_0-\Omega+(1+2x)^2\Psi-2xG$. Then:
\begin{enumerate}
\item[(a)] $E=R_1-(1+x)(\Omega+xH)$;
\item[(b)] for the planted tree, $\hat E=x\bigl[E+\Omega+2G+x(3+2x)H\bigr]$;
\item[(c)] for a merge of two rooted trees,
$E=E_1\Psi_2+E_2\Psi_1+x(H_1\Omega_2+H_2\Omega_1)-2G_1G_2$;
\item[(d)] for the merge of $q\ge1$ rooted trees with states indexed by $i$,
with $P_i=\prod_{l\ne i}\Psi_l$, $P_{il}=\prod_{t\ne i,l}\Psi_t$ and
$X_{il}=x(H_i\Omega_l+H_l\Omega_i)-2G_iG_l$:
$\Psi=\prod_i\Psi_i$, $X=\sum_iX_iP_i$ for $X\in\{H,G,\Omega\}$, and
$E=\sum_iE_iP_i+\sum_{i<l}X_{il}P_{il}$.
\end{enumerate}
\end{proposition}

\begin{proof}
(a) is Proposition~\ref{prop:rules}(b) rewritten:
$R_1=E+\Omega+x(1+x)H+x\Omega$.

(b) By Proposition~\ref{prop:rules}(c) and (a),
$\hat E=\hat R_0-\hat\Omega+(1+2x)^2\hat\Psi-2x\hat G$ with
$\hat R_0=(1+x)R_0+x\Omega$, $\hat\Omega=R_1=E+(1+x)(\Omega+xH)$,
$\hat\Psi=\Phi_1=(1+x)\Psi+xH$, $\hat G=x(xH+G)$, and
$R_0=E+\Omega-(1+2x)^2\Psi+2xG$. Substituting, the terms in $\Psi$ cancel
($-(1+x)(1+2x)^2+(1+2x)^2(1+x)$), and
\[
\hat E=xE+x\Omega+\bigl[2x(1+x)-2x^2\bigr]G+\bigl[x(1+2x)^2-x(1+x)-2x^3\bigr]H ,
\]
with $x(1+2x)^2-x(1+x)-2x^3=3x^2+2x^3$.

(c) By Proposition~\ref{prop:rules}(d),
$E=[R_{0,1}-\Omega_1-2xG_1]\Psi_2+[R_{0,2}-\Omega_2-2xG_2]\Psi_1+2(1+2x)^2\Psi_1\Psi_2+x(H_1\Omega_2+H_2\Omega_1)-2G_1G_2$,
and $(1+2x)^2\Psi_1\Psi_2$ is added once to each bracket.

(d) Induction on $q$, merging the first $q$ trees (state with index $0$) with
the last by (c): with $\Psi_0=\prod_{i\le q}\Psi_i$ and
$X_0=\sum_{i\le q}X_i\prod_{l\le q,l\ne i}\Psi_l$, the terms
$x(H_0\Omega_{q+1}+H_{q+1}\Omega_0)-2G_0G_{q+1}$ equal
$\sum_{i\le q}X_{i,q+1}\prod_{l\le q,l\ne i}\Psi_l$, and $E_0\Psi_{q+1}$
supplies the terms with both indices at most $q$.
\end{proof}

\begin{proposition}[two planted trees]\label{prop:S35}
Let $(A_1,u_1)$, $(A_2,u_2)$ be rooted trees of orders $a_1,a_2\ge1$, with the
quantities of Proposition~\ref{prop:S33} indexed by $1,2$, and let a hat
denote the state of the planted tree $(\hat A_i,r_i)$. Then
\[
x(\hat H_1\hat\Omega_2+\hat H_2\hat\Omega_1)-2\hat G_1\hat G_2
=x(\Psi_1R_{1,2}+\Psi_2R_{1,1})+x^2(H_1Q^1_2+H_2Q^1_1)
+x^4\sum_{v\in A_1\setminus u_1}\ \sum_{v'\in A_2\setminus u_2}(\delta_v-\delta_{v'})^2a_va_{v'},
\]
where $R_{1,i}=R_1(A_i,u_i)$ and each distance is taken in its own tree from
its own root. If $R_{1,i}\in\K_{a_i+1}$ for $i=1,2$, the left side is in
$\K_{a_1+a_2+2}$.
\end{proposition}

\begin{proof}
By Proposition~\ref{prop:rules}(c), $\hat H_i=\Phi_{0,i}=\Psi_i+xH_i$,
$\hat\Omega_i=R_{1,i}$ and $\hat G_i=x(xH_i+G_i)=x^2J_i$
(Proposition~\ref{prop:S33}(c)). By Proposition~\ref{prop:S33}(e),
$R_{1,i}=Q^1_i+x^2W_i$. So $x\hat H_1\hat\Omega_2=x\Psi_1R_{1,2}+x^2H_1Q^1_2+x^4H_1W_2$,
and the left side is
$x(\Psi_1R_{1,2}+\Psi_2R_{1,1})+x^2(H_1Q^1_2+H_2Q^1_1)+x^4\bigl[H_1W_2+H_2W_1-2J_1J_2\bigr]$.
With $H_i=\sum a_v$, $J_i=\sum\delta_va_v$, $W_i=\sum\delta_v^2a_v$, the bracket is
$\sum_{v,v'}a_va_{v'}(\delta_{v'}^2+\delta_v^2-2\delta_v\delta_{v'})$.

Membership: $\Psi_i\in\K_{a_i-1}$, so $x\Psi_1R_{1,2}\in\K_{2+(a_1-1)+(a_2+1)}$;
$H_1\in\K_{a_1-2}$ and $Q^1_2\in\K_{a_2}$, so $x^2H_1Q^1_2\in\K_{a_1+a_2+2}$;
$a_v\in\K_{a_1-2}$, $a_{v'}\in\K_{a_2-2}$, so
$x^4a_va_{v'}\in\K_{a_1+a_2+4}$. (If $a_i=1$ then $H_i=G_i=0$ and the sums
over $A_i\setminus u_i$ are empty.)
\end{proof}

\section{Proof of the lower bound}\label{sec:lowerproof}

\begin{theorem}\label{thm:T5}
For every rooted tree $(A,u)$ of order $a\ge1$:
\textup{(i)} $\Omega\in\K_a$; \textup{(ii)} $E\in\K_{a+1}$;
\textup{(iii)} $R_1(A,u)\in\K_{a+1}$.
\end{theorem}

\begin{proof}
(iii) follows from (i) and (ii): by Proposition~\ref{prop:S34}(a),
$R_1=E+(1+x)\Omega+x(1+x)H$, with $(1+x)\Omega\in\K_{a+1}$ and
$x(1+x)H\in\K_{(a-2)+3}$ (Proposition~\ref{prop:S33}(b), (K2)).

(i) and (ii) by induction on $a$. For $a=1$, $\Omega=0$ and
$E=-(1+2x)^2-0+(1+2x)^2-0=0$ (Proposition~\ref{prop:rules}(a)). Let $a\ge2$,
with branches $(A_i,c_i)$ of orders $a_i$, $\sum_ia_i=a-1$; (i)--(iii) hold for
each branch by induction, and $(A,u)$ is the merge of the planted trees
$(\hat A_i,r_i)$, whose $\hat\Psi_i=\Phi_1(A_i,c_i)$ lie in $\K_{a_i}$ by
Lemma~\ref{lem:K}.

(i) $\Omega=\sum_iR_1(A_i,c_i)\prod_{l\ne i}\Phi_1(A_l,c_l)$, and each term is
in $\K_{(a_i+1)+\sum_{l\ne i}a_l}=\K_a$.

(ii) By Proposition~\ref{prop:S34}(d) for the planted trees,
$E=\sum_i\hat E_i\hat P_i+\sum_{i<l}\hat X_{il}\hat P_{il}$, where the $\hat P$
are products of the $\hat\Psi_t\in\K_{a_t}$. By Proposition~\ref{prop:S34}(b),
$\hat E_i=x[E_i+\Omega_i+2G_i+x(3+2x)H_i]$ with the state of $(A_i,c_i)$; by
induction and Proposition~\ref{prop:S33}, $xE_i\in\K_{a_i+3}$,
$x\Omega_i\in\K_{a_i+2}$, $2xG_i=2x^2(G_i/x)\in\K_{a_i+3}$, $x^2H_i\in\K_{a_i+2}$
and $x^3H_i\in\K_{a_i+4}$; so $\hat E_i\in\K_{a_i+2}$ and
$\hat E_i\hat P_i\in\K_{a+1}$. By (iii) for the branches and
Proposition~\ref{prop:S35}, $\hat X_{il}\in\K_{a_i+a_l+2}$, so
$\hat X_{il}\hat P_{il}\in\K_{a+1}$.
\end{proof}

\begin{remark}[where the difficulty goes]\label{rem:difficulty}
In the merge rule the only negative term is $-2G_1G_2$. For planted trees it
is $-2x^4J_1J_2$, a product of first moments of the distances to the roots;
with the zeroth moments $H_i$ and the second moments $x^2W_i$ contained in
$R_{1,i}$ it forms the sum of squares of Proposition~\ref{prop:S35}. What
remains of $R_{1,i}$ is $Q^1_i$, one level lower in the cones, but it is
multiplied by $x^2H_l$, which carries one factor $x$ more than the term
$x\Psi_l$ that multiplies all of $R_{1,i}$. Direct bounds on the closed form
of Proposition~\ref{prop:S33}(e) fail by exactly one level: the edge terms of
$Q^1$ lie in $\K_a$, not $\K_{a+1}$.
\end{remark}

\begin{theorem}[the lower bound]\label{thm:lower}
\textup{(a)} For every tree $T$ of order $n\ge1$ and every vertex $u$,
$r_u(T)-(1+4x)(1+x)^{n-1}\in\K_n$.

\textup{(b)} For every tree $T$ of order $n\ge2$,
$\ell_T-\ell_{S_n}\in\K_{n-1}$; in particular $d_k(T)\ge d_k(S_n)$ for
$0\le k\le n-2$.

\textup{(c)} For every tree $T$ of order $n\ge3$, with $m=\lfloor n/2\rfloor$,
\[
d_m(T)-d_{m-1}(T)\ \ge\ d_m(S_n)-d_{m-1}(S_n)\ >\ 0 ,
\]
and more generally $d_\mu(T)-d_{\mu-1}(T)\ge d_\mu(S_n)-d_{\mu-1}(S_n)>0$
for $1\le\mu\le m$. Hence $d_0(T)<d_1(T)<\dots<d_m(T)$, and every maximizing
index of $d_0(T),\dots,d_{n-2}(T)$ is at least $\lfloor n/2\rfloor$.
\end{theorem}

\begin{proof}
(a) $\Psi_u=\prod_i\Phi_1(T_i,c_i)$ over the branches, with
$\sum_i|T_i|=n-1$. Write $\Phi_1(T_i,c_i)=(1+x)^{|T_i|}+b_i$ with
$b_i\in\K_{|T_i|}$ (Corollary~\ref{cor:phi1}) and expand: $\Psi_u-(1+x)^{n-1}$
is a sum of products with one factor from each $\K_{|T_i|}$, so it lies in
$\K_{n-1}$, and $(1+4x)[\Psi_u-(1+x)^{n-1}]\in\K_n$ by (K2). And
$\Omega_u\in\K_n$ by Theorem~\ref{thm:T5}(i).

(b) Build $T$ leaf by leaf as in Corollary~\ref{cor:leaf}(c), and the star
from its centre. Subtracting the two instances of Corollary~\ref{cor:leaf}(c)
(the star has $r_c(T_k)=(1+4x)(1+x)^{k-1}$, proof of
Corollary~\ref{cor:leaf}(d)):
\[
\ell_T-\ell_{S_n}=\sum_{k=1}^{n-1}(1+x)^{n-1-k}\bigl[r_{u_k}(T_k)-(1+4x)(1+x)^{k-1}\bigr],
\]
and by (a) for $T_k$ each bracket is in $\K_k$, so each term is in
$\K_{n-1}$. The last assertion follows from (K1) and Lemma~\ref{lem:ident}.

(c) Let $1\le\mu\le m$. Then $2\mu-1\le n-1$, so $\K_{n-1}\subseteq\K_{2\mu-1}$,
and $\lambda_\mu(\ell_T-\ell_{S_n})\ge0$ by (b) and Lemma~\ref{lem:A}. As
$n\ge3$, $\mu\le m\le n-2$, so $\lambda_\mu(\ell_T)=d_\mu(T)-d_{\mu-1}(T)$ by
Lemma~\ref{lem:ident}, and likewise for the star. That
$d_\mu(S_n)-d_{\mu-1}(S_n)>0$ is Lemma~\ref{lem:starshape}(a) below (for
$\mu=m$, Lemma~\ref{lem:starvalue}). So $d_0(T)<\dots<d_m(T)$, and no index
below $m$ is maximizing.
\end{proof}

\begin{lemma}\label{lem:starvalue}
Let $n\ge3$, $m=\lfloor n/2\rfloor$, $N=n-2$. Then
$d_m(S_n)-d_{m-1}(S_n)>0$; explicitly, $=4$ for $n=3$;
$=\binom{2m-2}{m-2}\frac{(2m-1)(3m-5)}{(m-1)(m+1)}$ for $n=2m\ge4$; and
$=4\binom{2m-1}{m-2}\bigl[\frac{2(2m-1)}{m-1}-\frac4{m+2}\bigr]$ for $n=2m+1\ge5$.
\end{lemma}

\begin{proof}
By Corollary~\ref{cor:leaf}(d),
$[x^k]\ell_{S_n}=(n-1)\binom Nk+4(n-2)\binom N{k-1}-4\binom N{k-2}$, so
\[
\lambda_m(\ell_{S_n})=(n-1)\Bigl[\tbinom Nm-\tbinom N{m-1}\Bigr]+4(n-2)\Bigl[\tbinom N{m-1}-\tbinom N{m-2}\Bigr]-4\Bigl[\tbinom N{m-2}-\tbinom N{m-3}\Bigr].
\]
For $n=3$: $\ell=(1+x)(2+4x-4x^2)$, $d_0=2$, $d_1=6$. For $n=2m\ge4$, with
$\beta=\binom{2m-2}{m-2}=\binom{2m-2}{m}$, $\binom{2m-2}{m-1}=\beta\frac m{m-1}$,
$\binom{2m-2}{m-3}=\beta\frac{m-2}{m+1}$:
$\lambda_m=\beta\bigl[\frac{3n-7}{m-1}-\frac{12}{m+1}\bigr]$, and
$(6m-7)(m+1)-12(m-1)=(2m-1)(3m-5)>0$. For $n=2m+1\ge5$:
$\binom{2m-1}m=\binom{2m-1}{m-1}$, and with $\beta=\binom{2m-1}{m-2}$,
$\binom{2m-1}{m-1}=\beta\frac{m+1}{m-1}$, $\binom{2m-1}{m-3}=\beta\frac{m-2}{m+2}$:
$\lambda_m=4\beta\bigl[\frac{2(2m-1)}{m-1}-\frac4{m+2}\bigr]>0$ since
$\frac{2(2m-1)}{m-1}>4$.
\end{proof}

\begin{lemma}[the star]\label{lem:starshape}
Let $n\ge3$ and $m=\lfloor n/2\rfloor$. \textup{(a)}
$d_0(S_n)<d_1(S_n)<\dots<d_m(S_n)$. \textup{(b)}
$d_m(S_n)>d_{m+1}(S_n)>\dots>d_{n-2}(S_n)$. So $\lfloor n/2\rfloor$ is the
unique maximizing index of the star.
\end{lemma}

\begin{proof}
Let $N=n-2$ and $F(j)=\binom Nj-\binom N{j-1}$, so that $F(j)=0$ for $j<0$
and $F(N+1-j)=-F(j)$ for all $j$. As in the proof of
Lemma~\ref{lem:starvalue},
\[
\lambda_k(\ell_{S_n})=(n-1)F(k)+4(n-2)F(k-1)-4F(k-2)\qquad(1\le k\le n-2).
\]
For $0\le j\le N/2$, $F(j)=\binom Nj\frac{N-2j+1}{N-j+1}>0$; for $0\le j<(N-1)/2$
(so that $F(j+1)>0$),
\begin{equation}\label{eq:Fratio}
\frac{F(j)}{F(j+1)}=\frac{(j+1)(N-2j+1)}{(N-j+1)(N-2j-1)} .
\end{equation}

(a) For $k=m$ this is Lemma~\ref{lem:starvalue}. Let $1\le k\le m-1$, so
$2k\le N$, and $F(k)>0$, $F(k-1)>0$. If $k=1$, $F(-1)=0$ and $\lambda_1>0$. If
$k\ge2$, apply~\eqref{eq:Fratio} with $j=k-2$: then $2j<N$ and
$N-2j-1=N-2k+3\ge3$, so $F(k-2)/F(k-1)<1+2/3=5/3$. As $n\ge2k+2\ge6$,
$4F(k-2)<\frac{20}3F(k-1)<4(n-2)F(k-1)$, and $\lambda_k>(n-1)F(k)>0$.

(b) Let $m+1\le k\le n-2$ (nothing to prove unless $n\ge5$). Then
$k>(N+1)/2$, so $F(k)<0$, and $k-1\ge m\ge(N+1)/2$, so $F(k-1)\le0$. If
$F(k-2)\ge0$, then $\lambda_k\le(n-1)F(k)<0$. Otherwise put $j=N+2-k$; then
$j\ge2$ (as $k\le N$), $j<(N-1)/2$ (since $F(k-2)<0$ forces $k-2>(N+1)/2$),
$-F(k-1)=F(j)>0$
and $-F(k-2)=F(j+1)>0$, and by~\eqref{eq:Fratio}
$F(j+1)/F(j)<(N-j+1)/(j+1)\le(N-1)/3<n-2$. So
$-4F(k-2)<4(n-2)(-F(k-1))$, and
$\lambda_k<(n-1)F(k)+4(n-2)F(k-1)-4(n-2)F(k-1)=(n-1)F(k)<0$.
\end{proof}

\begin{remark}\label{rem:star-peak}
By Lemma~\ref{lem:starshape} the lower bound is attained by the star for every
$n\ge3$, in agreement with~\cite[Theorem~1]{Col89}. For even $n$ the relative
margin $1-d_{m-1}/d_m$ of the star is of order $1/n$; this is the case in
which the lower bound is delicate.
\end{remark}

\begin{remark}[on the use of log-concavity]\label{rem:noLC}
Theorem~\ref{thm:lower}(c) does not use (LC): the cone comparison holds for
every increment $\lambda_\mu$ with $\mu\le m$, and the star increases strictly
up to $m$ (Lemma~\ref{lem:starshape}(a)).
\end{remark}

\begin{remark}[a sharper cone]\label{rem:Ccone}
Let $\mathcal C_d\subseteq\K_d$ be the cone spanned by the $x^e(1+x)^i$ with
$e+i\le d\le2e+i$. A polynomial of $\K_d$ of degree at most $d$ lies in
$\mathcal C_d$ (no generator of positive weight can have degree above $d$).
Since $\ell_T-\ell_{S_n}$ has degree at most $n-2$ and $R_1(A,u)$ degree at
most $a+1$, Theorem~\ref{thm:lower}(b) and Theorem~\ref{thm:T5}(iii) also hold
with $\mathcal C$ in place of $\K$. This is not needed here.
\end{remark}
\part*{The upper bound}

\section{The spectral form of \texorpdfstring{$\ell_T$}{l\_T}}\label{sec:spectral}

Notation for this part. For a graph $H$ write $G_H=(I+L_H)^{-1}$ and
$\Tr(H)=\tr G_H$. For a rooted tree $(X,v)$ put
\[
g_X=(G_X)_{vv},\quad k_X=(G_X^2)_{vv},\quad \varepsilon_X=1/g_X-1,\quad
\kappa_X=k_X/g_X^2,\quad t_X=\tr G_X .
\]
The letters $q$, $\beta$, $G$, $E$, $a$ have in this part the meanings given
here, unrelated to those of Sections~\ref{sec:rooted}--\ref{sec:lowerproof}.

\begin{lemma}\label{lem:S40}
Let $T$ be a tree of order $n\ge2$. The matrix $D/2$ has one positive
eigenvalue $\alpha$ and $n-1$ negative eigenvalues
$-\beta_1,\dots,-\beta_{n-1}$, and $\alpha=\sum_i\beta_i$. With
$q(x)=\prod_{i=1}^{n-1}(x+\beta_i)=\sum_{k=0}^{n-1}q_kx^k$,
\[
\ell_T(x)=4(\alpha-x)q(x),\qquad\text{so}\qquad d_k=4(\alpha q_k-q_{k-1})
\quad(0\le k\le n-2,\ q_{-1}=0).
\]
\end{lemma}

\begin{proof}
By (GL), $2D^{-1}=-L+ww^{\mathsf T}$ with $w=\tau/\sqrt{n-1}$. Since a tree is
connected, the eigenvalues of $-L$ in decreasing order are
$0\ge-\mu_{(1)}\ge\dots\ge-\mu_{(n-1)}$, where
$0<\mu_{(1)}\le\dots\le\mu_{(n-1)}$ are the non-zero Laplacian eigenvalues;
thus $\lambda_1(-L)=0$ and $\lambda_i(-L)=-\mu_{(i-1)}$ for $i\ge2$. Let
$\nu_1\ge\dots\ge\nu_n$ be the eigenvalues of $2D^{-1}$. By (W),
$\nu_i\le\lambda_{i-1}(-L)\le0$ for $i\ge2$; none is zero, since $D$ is
invertible by (GP); so $\nu_2,\dots,\nu_n<0$. The product $\prod_i\nu_i=\det(2D^{-1})$
has the sign of $\det D$, that is $(-1)^{n-1}$ by (GP); hence $\nu_1>0$. The
eigenvalues of $D/2$ are the $1/\nu_i$: $\alpha=1/\nu_1>0$ and, for $i\ge2$,
$-\beta_{i-1}=1/\nu_i<0$. Since $\tr D=0$, $\alpha=\sum_i\beta_i$. Finally,
with $\theta_1,\dots,\theta_n$ the eigenvalues of $D/2$,
$\det(2xI-D)=2^n\prod_i(x-\theta_i)=2^n(x-\alpha)q(x)$, so
$\ell_T=-2^{-(n-2)}2^n(x-\alpha)q(x)=4(\alpha-x)q(x)$. Comparing coefficients
and using Lemma~\ref{lem:ident} gives $d_k=4(\alpha q_k-q_{k-1})$ for
$0\le k\le n-2$.
\end{proof}

\begin{lemma}[the peak of $d$ lies at or before a descent of $q$]\label{lem:S41}
Let $T$ be a tree of order $n\ge3$, with $\alpha$, $q$ as in
Lemma~\ref{lem:S40}, and let $0\le k\le n-3$. If $q_{k+1}<q_k$, then
$d_{j+1}<d_j$ for every $j$ with $k\le j\le n-3$; hence every maximizing
index of $d_0,\dots,d_{n-2}$ is at most $k$.
\end{lemma}

\begin{proof}
The polynomial $q$ has degree $n-1$, positive coefficients (its zeros
$-\beta_i$ are negative) and only real zeros. By (N),
$q_j^2>q_{j-1}q_{j+1}$ for $1\le j\le n-2$, so the ratios
$\varrho_j=q_{j+1}/q_j$, $0\le j\le n-2$, are strictly decreasing. Since $q$
is monic, $q_{n-2}=\sum_i\beta_i=\alpha$, so $\varrho_{n-2}=1/\alpha$, and
hence $\alpha\varrho_j>1$ for $0\le j\le n-3$. Fix $j$ with $k\le j\le n-3$;
then $\varrho_j\le\varrho_k<1$. By Lemma~\ref{lem:S40},
\[
d_{j+1}-d_j=4\bigl[\alpha(q_{j+1}-q_j)-(q_j-q_{j-1})\bigr].
\]
If $j=0$, this is $4[\alpha(q_1-q_0)-q_0]<0$. If $j\ge1$ and
$q_j\ge q_{j-1}$, the first term is negative and the second is non-positive.
If $j\ge1$ and $q_j<q_{j-1}$, then $\varrho_j<\varrho_{j-1}<1$, and writing
$q_j=\varrho_{j-1}q_{j-1}$, $q_{j+1}=\varrho_j\varrho_{j-1}q_{j-1}$,
\[
d_{j+1}-d_j=4q_{j-1}\bigl[(1-\varrho_{j-1})-\alpha\varrho_{j-1}(1-\varrho_j)\bigr];
\]
since $\alpha\varrho_{j-1}>1$ and $1-\varrho_j>1-\varrho_{j-1}>0$, the bracket
is negative. So $d_k>d_{k+1}>\dots>d_{n-2}$, and no index larger than $k$ is
maximizing.
\end{proof}

\begin{lemma}[the mean of $q$ and the Laplacian spectrum]\label{lem:S42}
Let $T$ be a tree of order $n\ge2$, with $q$ as in Lemma~\ref{lem:S40}. Then
\[
\mu_q:=\frac{q'(1)}{q(1)}=\sum_{i=1}^{n-1}\frac1{1+\beta_i}
\le\sum_{j=1}^{n-1}\frac{\mu_{(j)}}{1+\mu_{(j)}}=n-\Tr(T).
\]
\end{lemma}

\begin{proof}
The first equality is the logarithmic derivative of $\prod_i(x+\beta_i)$ at
$x=1$. In the notation of the proof of Lemma~\ref{lem:S40},
$\beta_{i-1}=-1/\nu_i=1/|\nu_i|$ for $2\le i\le n$, so
$1/(1+\beta_{i-1})=|\nu_i|/(1+|\nu_i|)$. By (W), $\nu_i\ge\lambda_i(-L)=-\mu_{(i-1)}$
for $i\ge2$, so $|\nu_i|\le\mu_{(i-1)}$; and $y\mapsto y/(1+y)$ is increasing
on $(0,\infty)$. This gives the inequality. Finally
$\Tr(T)=1+\sum_j1/(1+\mu_{(j)})$ (the eigenvalue $0$ of $L$ contributes $1$),
so $\sum_j\mu_{(j)}/(1+\mu_{(j)})=\sum_j\bigl(1-1/(1+\mu_{(j)})\bigr)=(n-1)-(\Tr(T)-1)=n-\Tr(T)$.
\end{proof}

\begin{remark}
The interlacing between the distance eigenvalues and the reciprocals of the
non-zero Laplacian eigenvalues used in Lemmas~\ref{lem:S40}
and~\ref{lem:S42} is a theorem of Merris~\cite{Mer90}:
$0>-2/\mu_{(n-1)}\ge\partial_2\ge-2/\mu_{(n-2)}\ge\partial_3\ge\dots\ge-2/\mu_{(1)}\ge\partial_n$
(revision~2: the last inequality was written in the wrong direction in the
first version),
where $\partial_1\ge\dots\ge\partial_n$ are the eigenvalues of $D$; the
inertia of $D$ is $(1,n-1,0)$. Merris's paper was read in full by the
project's literature audit; the equation number cited for this chain in the
project notes was not confirmed by a reviewer. The proofs above are the
same argument, through (GL).
\end{remark}

\section{The trace of \texorpdfstring{$(I+L)^{-1}$}{(I+L)\^{}(-1)}}\label{sec:trace}

\begin{lemma}[gluing two rooted trees]\label{lem:S43}
Let $(X,v)$ and $(Y,v)$ be rooted trees and let $T$ be the tree obtained by
identifying their roots. Then $1/g_T=1/g_X+1/g_Y-1$ (so
$\varepsilon_T=\varepsilon_X+\varepsilon_Y$), and
\[
t_T=t_X+t_Y-g_T\bigl[1+\varepsilon_Y\kappa_Xg_X+\varepsilon_X\kappa_Yg_Y\bigr].
\]
\end{lemma}

\begin{proof}
If $Y=\{v\}$ then $g_Y=1$, $\varepsilon_Y=0$, $t_Y=1$, $T=X$, and both claims
are immediate; similarly if $X=\{v\}$. Otherwise order the vertices of $T$ as
the vertices of $X$ (including $v$), then those of $Y\setminus v$. Let
$C=(I+L_Y)[Y\setminus v]$ be the principal submatrix of $I+L_Y$ on
$Y\setminus v$ (positive definite, as a principal submatrix of the positive
definite matrix $I+L_Y$), let $b\in\mathbb R^{Y\setminus v}$ be the indicator
vector of the neighbours of $v$ in $Y$, and let $\tilde b$ be the
$|X|\times|Y\setminus v|$ matrix whose row $v$ is $b^{\mathsf T}$ and whose
other rows are zero. Then
\[
I+L_T=\begin{pmatrix}I+L_X+\deg_Y(v)E_{vv}&-\tilde b\\-\tilde b^{\mathsf T}&C\end{pmatrix}.
\]
By the Schur complement formula, the $X$-block of $G_T$ is
$\bigl(I+L_X+sE_{vv}\bigr)^{-1}$ with $s=\deg_Y(v)-b^{\mathsf T}C^{-1}b$
(because $\tilde bC^{-1}\tilde b^{\mathsf T}=(b^{\mathsf T}C^{-1}b)E_{vv}$).
In $Y$ alone, ordered with $v$ first,
$I+L_Y=\begin{pmatrix}1+\deg_Y(v)&-b^{\mathsf T}\\-b&C\end{pmatrix}$, and the
same elimination gives $1/g_Y=1+\deg_Y(v)-b^{\mathsf T}C^{-1}b$; hence
$s=\varepsilon_Y$. By the Sherman--Morrison formula,
\[
(I+L_X+sE_{vv})^{-1}=G_X-\frac{s\,G_Xe_ve_v^{\mathsf T}G_X}{1+s\,g_X};
\]
its $(v,v)$ entry is $g_X/(1+sg_X)$, and its trace is
$t_X-sk_X/(1+sg_X)$ (as $\tr(G_Xe_ve_v^{\mathsf T}G_X)=e_v^{\mathsf T}G_X^2e_v=k_X$).
Hence $g_T=g_X/(1+\varepsilon_Yg_X)$, that is $1/g_T=1/g_X+\varepsilon_Y$,
which is the first claim. Exchanging the roles of $X$ and $Y$, the
$Y$-block of $G_T$ (on the vertex set of $Y$, including $v$) has trace
$t_Y-\varepsilon_Xk_Y/(1+\varepsilon_Xg_Y)$. The two blocks are principal
submatrices of $G_T$ whose index sets share only $v$, so
\[
t_T=t_X-\frac{\varepsilon_Yk_X}{1+\varepsilon_Yg_X}+t_Y-\frac{\varepsilon_Xk_Y}{1+\varepsilon_Xg_Y}-g_T .
\]
Finally
$\varepsilon_Yk_X/(1+\varepsilon_Yg_X)=\varepsilon_Y\kappa_Xg_X\cdot g_X/(1+\varepsilon_Yg_X)=g_T\,\varepsilon_Y\kappa_Xg_X$,
and symmetrically $\varepsilon_Xk_Y/(1+\varepsilon_Xg_Y)=g_T\,\varepsilon_X\kappa_Yg_Y$.
\end{proof}

\begin{lemma}[bounds at a root]\label{lem:S43p}
Let $(X,v)$ be a rooted tree with $|X|\ge2$. Then $0<g_X<1$, so
$a:=1/g_X=1+\varepsilon_X>1$, and $1\le\kappa_X\le1+\varepsilon_X/2=(1+a)/2$.
\end{lemma}

\begin{proof}
Write $G=G_X$. By (M), row $v$ of $G$ is positive and sums to $1$; as
$|X|\ge2$, $0<G_{vv}<1$. Hence
$\varepsilon_X=(1-G_{vv})/G_{vv}=\sum_{w\ne v}G_{vw}/G_{vv}>0$, and, $G$ being
symmetric, $\kappa_X=\sum_wG_{vw}^2/G_{vv}^2=1+\sum_{w\ne v}(G_{vw}/G_{vv})^2\ge1$.
For the upper bound it is enough to show $G_{wv}\le G_{vv}/2$ for every
$w\ne v$, since then
$\sum_{w\ne v}(G_{vw}/G_{vv})^2\le\frac12\sum_{w\ne v}G_{vw}/G_{vv}=\varepsilon_X/2$.
Let $w^*$ maximize $G_{wv}$ over $w\ne v$, and let $d=\deg w^*\ge1$. Row
$w^*$ of $(I+L)G=I$, in column $v$, reads $(1+d)G_{w^*v}=\sum_{x\sim w^*}G_{xv}$.
If $v$ is not a neighbour of $w^*$, the right side is at most $dG_{w^*v}$,
so $G_{w^*v}\le0$, contradicting (M). So $v\sim w^*$, and
$(1+d)G_{w^*v}\le G_{vv}+(d-1)G_{w^*v}$, that is $2G_{w^*v}\le G_{vv}$.
\end{proof}

For $N\ge2$ let $P=P_{N+1}$ with vertices $0,1,\dots,N$ in order, $G=G_P$,
and for a vertex $r$ put $E_r=1/G_{rr}-1$ and $c_r=(G^2)_{rr}/G_{rr}$.

\begin{lemma}[the path seen from one vertex]\label{lem:S44}
Let $N\ge2$ and $1\le r\le N-1$. Then (i) $E_r>E_0$; (ii) $c_r>3/5$;
(iii) $c_0\ge c_r$, with equality only for $N=2$, $r=1$.
\end{lemma}

\begin{proof}
Let $G(t)=(tI+L_P)^{-1}$ for $t>0$, and write $t=4\sinh^2u$ with $u>0$. Then
$t=1$ corresponds to $u=\log\varphi$, since
$\sinh(\log\varphi)=(\varphi-\varphi^{-1})/2=1/2$. Let $e_j(t)$ be the
determinant of the $j\times j$ tridiagonal matrix with diagonal
$(t+1,t+2,\dots,t+2)$ and entries $-1$ next to the diagonal ($e_0=1$); it is
the principal minor of $tI+L_P$ on a segment of $j$ consecutive vertices
that contains an end of $P$ and does not contain $r$. Expanding along the
last row, $e_j=(t+2)e_{j-1}-e_{j-2}$ for $j\ge2$. Since $t+2=2\cosh2u$ and
$e_1=t+1=2\cosh2u-1=\cosh3u/\cosh u$, and since
$2\cosh2u\cosh((2j-1)u)=\cosh((2j+1)u)+\cosh((2j-3)u)$, induction gives
\[
e_j=\cosh((2j+1)u)/\cosh u\qquad(j\ge0).
\]
For $m\ge2$, expanding $\det(tI+L_{P_m})$ at an end vertex gives
$(t+1)e_{m-1}-e_{m-2}$; with $t+1=2\cosh2u-1$ and
$2\cosh2u\cosh((2m-1)u)=\cosh((2m+1)u)+\cosh((2m-3)u)$ this equals
\[
\frac{\cosh((2m+1)u)-\cosh((2m-1)u)}{\cosh u}=\frac{2\sinh u\,\sinh(2mu)}{\cosh u};
\]
for $m=1$ the right side is $4\sinh^2u=t=\det(tI+L_{P_1})$. Removing $r$
splits $P$ into segments of $r$ and $s=N-r$ vertices. With $a=2r+1$,
$b=2s+1$, $a+b=2N+2$, Cramer's rule gives, for every $0\le r\le N$,
\begin{equation}\label{eq:Grr}
G_{rr}(t)=\frac{e_re_s}{\det(tI+L_P)}=\frac{\cosh(au)\cosh(bu)}{2\sinh u\cosh u\,\sinh((a+b)u)}.
\end{equation}

(i) By \eqref{eq:Grr}, for $r=0$ one has $a=1$ and $b=a+b-1$ with the same
$a+b$; so
\[
\frac{G_{rr}}{G_{00}}=\frac{\cosh(au)\cosh(bu)}{\cosh(u)\cosh((a+b-1)u)}
=\frac{\cosh((a+b)u)+\cosh((a-b)u)}{\cosh((a+b)u)+\cosh((a+b-2)u)}<1,
\]
because $|a-b|<a+b-2$ when $a,b\ge3$, that is when $1\le r\le N-1$. So
$G_{rr}<G_{00}$ at $t=1$, and $E_r>E_0$.

(ii) Since $\frac{d}{dt}G(t)=-G(t)^2$, $c_r=-\frac{d}{dt}\log G_{rr}(t)$ at
$t=1$. From \eqref{eq:Grr},
\[
\frac{d}{du}\log G_{rr}=a\tanh(au)+b\tanh(bu)-\tanh u-\coth u-(a+b)\coth((a+b)u),
\]
and $dt/du=4\sinh2u$. At $u=\log\varphi$: $\cosh u=\sqrt5/2$,
$\sinh2u=(\varphi^2-\varphi^{-2})/2=\sqrt5/2$, $\tanh u=1/\sqrt5$,
$\coth u=\sqrt5$. Hence
\begin{equation}\label{eq:cr}
c_r=\frac{6/\sqrt5+(a+b)\coth((a+b)u)-a\tanh(au)-b\tanh(bu)}{2\sqrt5},
\end{equation}
valid for every $0\le r\le N$ (for $r=0$, $a=1$ and $b=2N+1$). Since
$\coth>1>\tanh$ on $(0,\infty)$, $(a+b)\coth((a+b)u)>a+b>a\tanh(au)+b\tanh(bu)$,
so the numerator exceeds $6/\sqrt5$ and $c_r>3/5$.

(iii) In \eqref{eq:cr} the $\coth$ terms are the same for $r$ and for $0$.
Using $\tanh z=1-2/(e^{2z}+1)$, and noting that the linear parts cancel
because $a+b=1+(2N+1)$, one gets at $u=\log\varphi$ (where
$e^{2xu}=\varphi^{2x}$)
\[
c_0-c_r=\frac{\psi(1)+\psi(2N+1)-\psi(a)-\psi(b)}{\sqrt5},\qquad
\psi(x)=\frac{x}{\varphi^{2x}+1}.
\]
Put $f_k=\psi(2k+1)$ and $\delta_k=f_k-f_{k+1}$. Then $\psi(a)=f_r$,
$\psi(b)=f_s$, and, as $N=s+r$,
\[
f_0+f_N-f_r-f_s=(f_0-f_r)-(f_s-f_{s+r})=\sum_{i=0}^{r-1}(\delta_i-\delta_{s+i}).
\]
It suffices to show that $\delta$ is non-increasing with $\delta_0=\delta_1$
and $\delta_k>\delta_{k+1}$ for $k\ge1$. Then every term is $\ge0$ (as
$s+i>i$). If $r\ge2$, the term $i=1$ is $\delta_1-\delta_{s+1}>0$; if $r=1$,
the only term is $\delta_0-\delta_s$, which vanishes exactly when $s=1$. So
equality holds only for $r=s=1$, that is $N=2$, $r=1$.

\emph{$\delta_0=\delta_1$.} For odd $x$,
$\varphi^{2x}+1=\varphi^x(\varphi^x+\varphi^{-x})=\sqrt5F_x\varphi^x$, where
$F_x$ is the Fibonacci number (Binet's formula). With $F_1=1$, $F_3=2$,
$F_5=5$,
\[
\psi(1)-2\psi(3)+\psi(5)=\frac{\varphi^{-1}-3\varphi^{-3}+\varphi^{-5}}{\sqrt5}
=\frac{\varphi^{-5}(\varphi^4-3\varphi^2+1)}{\sqrt5}=0,
\]
since $\varphi^4=3\varphi+2$ and $3\varphi^2=3\varphi+3$.

\emph{$\delta_k>\delta_{k+1}$ for $k\ge1$.} Since
$\delta_k-\delta_{k+1}=\psi(2k+1)-2\psi(2k+3)+\psi(2k+5)$, this follows from
strict convexity of $\psi$ on $[3,\infty)$. With $\gamma=2\log\varphi$ and
$\sigma(x)=1/(e^{\gamma x}+1)$ one has $\psi=x\sigma$,
$\sigma'=-\gamma\sigma(1-\sigma)$, and
\[
\psi''(x)=\gamma\sigma(1-\sigma)\bigl[\gamma x\tanh(\gamma x/2)-2\bigr],
\]
which is positive if and only if $z\tanh z>1$ for $z=x\log\varphi$. The
function $z\tanh z$ is increasing on $(0,\infty)$, so it suffices that
$z\tanh z>1$ at $z=3\log\varphi$. There
$\tanh(3\log\varphi)=(\varphi^6-1)/(\varphi^6+1)=(8\varphi+4)/(8\varphi+6)=2/\sqrt5$;
and $\tanh(\log\varphi)=(\varphi-\varphi^{-1})/(\varphi+\varphi^{-1})=1/\sqrt5$,
so $\log\varphi=\operatorname{artanh}(1/\sqrt5)>1/\sqrt5$ (as
$\operatorname{artanh}t>t$ for $0<t<1$). Hence
$3\log\varphi\,\tanh(3\log\varphi)>(3/\sqrt5)(2/\sqrt5)=6/5>1$.
\end{proof}

\begin{lemma}[merging two pendant paths lowers the trace]\label{lem:S45}
Let $(X,v)$ be a rooted tree with $|X|\ge2$ and let $j_1,j_2\ge1$. Let
$T(j_1,j_2)$ be the tree obtained from $X$ by attaching at $v$ two pendant
paths with $j_1$ and $j_2$ new vertices, and $T(j_1+j_2,0)$ the tree
obtained from $X$ by attaching at $v$ one pendant path with $j_1+j_2$ new
vertices. Then $\Tr(T(j_1,j_2))>\Tr(T(j_1+j_2,0))$.
\end{lemma}

\begin{proof}
Let $N=j_1+j_2\ge2$ and $P=P_{N+1}$ as before Lemma~\ref{lem:S44}. Then
$T(j_1,j_2)$ is $X$ glued at $v$ to $(P,j_1)$, and $T(N,0)$ is $X$ glued at
$v$ to $(P,0)$; note $1\le j_1\le N-1$. In Lemma~\ref{lem:S43} with
$(Y,v)=(P,r)$ one has $\varepsilon_Y=E_r$, $\kappa_Yg_Y=k_Y/g_Y=c_r$, and
$t_Y=\Tr(P)$, which does not depend on $r$. Put $a=1/g_X$ and
$\kappa=\kappa_X$; then $\varepsilon_X=a-1$ and $g_T=1/(a+E_r)$, so
Lemma~\ref{lem:S43} gives
\begin{gather*}
\Tr(T(N,0))=t_X+\Tr(P)-\Theta(E_0,c_0),\\
\Tr(T(j_1,j_2))=t_X+\Tr(P)-\Theta(E_{j_1},c_{j_1}),\\
\Theta(E,c)=\frac{1+E\kappa/a+(a-1)c}{a+E}.
\end{gather*}
Write $E=E_{j_1}$, $c=c_{j_1}$, $E'=E_0$, $c'=c_0$. Then
$\Tr(T(j_1,j_2))-\Tr(T(N,0))=\Theta(E',c')-\Theta(E,c)$. Multiplying
$\Theta(E',c')$ by $a(a+E')$ gives $a[1+(a-1)c']+\kappa E'$, and similarly
for $\Theta(E,c)$; expanding, the terms in $\kappa$ give $\kappa a(E'-E)$, and
\begin{gather*}
a(a+E)(a+E')\bigl[\Theta(E',c')-\Theta(E,c)\bigr]=a\,\mathcal N,\\
\mathcal N=(E-E')(1-\kappa)+(a-1)\bigl[(a+E)c'-(a+E')c\bigr].
\end{gather*}
By Lemma~\ref{lem:S44}(i), $E>E'$, and by Lemma~\ref{lem:S43p},
$1-\kappa\ge-(a-1)/2$; so
$\mathcal N\ge(a-1)\bigl[(a+E)c'-(a+E')c-(E-E')/2\bigr]$. By
Lemma~\ref{lem:S44}(iii), $c'\ge c$, so $(a+E)c'\ge(a+E)c$ and
\[
\mathcal N\ge(a-1)(E-E')(c-1/2),
\]
which is positive by Lemma~\ref{lem:S43p} ($a>1$), Lemma~\ref{lem:S44}(i)
and Lemma~\ref{lem:S44}(ii) ($c>3/5$). As $a(a+E)(a+E')>0$, the claim
follows.
\end{proof}

\begin{theorem}[the path minimizes $\tr(I+L)^{-1}$]\label{thm:T10}
For every tree $T$ of order $n\ge1$, $\tr(I+L_T)^{-1}\ge\tr(I+L_{P_n})^{-1}$,
with equality only for $T\cong P_n$. Equivalently, among trees of order $n$,
$\sum_j\mu_j/(1+\mu_j)$ over the Laplacian eigenvalues $\mu_j$ is maximized
exactly by the path.
\end{theorem}

\begin{proof}
By induction on the number of leaves. A tree with at most two leaves is a
path. Otherwise $T$ has a vertex of degree at least $3$. Root $T$ at a leaf
$\rho$ and let $v$ be a vertex of degree $\ge3$ at maximum distance from
$\rho$. Every branch at $v$ (component of $T-v$) not containing $\rho$
contains no vertex of degree $\ge3$, since each of its vertices is farther
from $\rho$ than $v$; so it is a path, and its vertex adjacent to $v$ has
degree at most $2$ in $T$, hence is an end of that path. There are
$\deg v-1\ge2$ such branches. Take two of them, with $j_1$ and $j_2$
vertices, and let $X$ be $T$ without them; then $v\in X$ and $\rho\in X$
with $\rho\ne v$, so $|X|\ge2$, and $T=T(j_1,j_2)$ in the notation of
Lemma~\ref{lem:S45}. By Lemma~\ref{lem:S45}, $T'=T(j_1+j_2,0)$ has
$\Tr(T')<\Tr(T)$. The tree $T'$ has order $n$ and one leaf fewer than $T$:
the two ends of the two pendant paths are replaced by one, and $v$ keeps
degree $\deg v-1\ge2$. By induction $\Tr(T')\ge\Tr(P_n)$, so
$\Tr(T)>\Tr(P_n)$. The equivalent form follows from
$\sum_j\mu_j/(1+\mu_j)=n-\Tr(T)$ (proof of Lemma~\ref{lem:S42}).
\end{proof}

\begin{remark}[Theorem~\ref{thm:T10} follows from~\cite{AMW}]\label{rem:AMW}
Theorem~\ref{thm:T10} is a consequence of a theorem of Andriantiana,
Razanajatovo Misanantenaina and Wagner on the average size of matchings; the
proof above is independent of theirs, and a reader willing to rely
on~\cite{AMW} may skip Lemmas~\ref{lem:S43}--\ref{lem:S45}. For a graph $G$
let $m_k(G)$ be the number of matchings of $G$ with $k$ edges ($m_0=1$) and
$\mathrm{avm}(G)=\sum_kk\,m_k(G)/\sum_km_k(G)$. Put $a=(5-\sqrt5)/10$ and
$b=(7\sqrt5-17)/10$. (Numbering and pages of the arXiv version of~\cite{AMW};
the journal version was not compared.) By \cite[Theorem~4.4]{AMW}, every
tree $T'$ of order $N$ that is not a path has $\mathrm{avm}(T')\le aN+b$; its
proof includes, for a vertex of degree~$3$, $35$ cases that are checked by
computer \cite[pp.~13--15]{AMW}. The strict form
$\mathrm{avm}(P_N)>aN+b\ge\mathrm{avm}(T')$ is stated at the end of that
proof \cite[p.~16]{AMW}; it also follows from \cite[Lemma~4.2,
eq.~(8)]{AMW},
$\mathrm{avm}(P_N)=aN+\frac{1-\sqrt5}{10}-\frac{N+1}{\sqrt5\,((-\varphi^2)^{N+1}-1)}$;
since $\frac{1-\sqrt5}{10}-b=\frac{9-4\sqrt5}5>0.011$ and the last term is
positive for even $N$ and of absolute value at most
$6/(\sqrt5(\varphi^{12}-1))<0.0084$ for odd $N\ge5$, one has
$\mathrm{avm}(P_N)>aN+b$ for $N\ge4$. So the path is the unique maximizer of
$\mathrm{avm}$ among trees of order $N$, for every $N$.

Let $T$ be a tree of order $n\ge2$ with degrees $k_v$, $\mathcal A_T$ its
adjacency matrix, $\mathcal B$ its $n\times(n-1)$ unsigned vertex--edge
incidence matrix, and $S(T)$ its subdivision (each edge replaced by a path
of length two), a tree of order $N=2n-1$. Its adjacency matrix
$\mathcal A_{S(T)}$ is
$\left(\begin{smallmatrix}0&\mathcal B\\\mathcal B^\T&0\end{smallmatrix}\right)$,
and $\mathcal B\mathcal B^\T=\diag(k_v)+\mathcal A_T$ equals $\Sigma L_T\Sigma$
for the diagonal sign matrix $\Sigma$ of the bipartition of $T$. As $L_T$ has
the simple eigenvalue $0$, $\mathcal B$ has rank $n-1$, and the spectrum of
$S(T)$ consists of
$\pm\sqrt{\mu_j}$ for the $n-1$ positive Laplacian eigenvalues
$\mu_1,\dots,\mu_{n-1}$ of $T$, together with exactly one eigenvalue $0$. A
forest has no cycles, so in the permutation expansion of
$\det(yI-\mathcal A_{S(T)})$ only products of disjoint transpositions along
edges survive, and $\det(yI-\mathcal A_{S(T)})=\sum_k(-1)^km_k(S(T))\,y^{N-2k}$. Comparing with
$y\prod_j(y^2-\mu_j)$ gives $\sum_km_k(S(T))\,t^k=\prod_{j=1}^{n-1}(1+\mu_jt)$,
and the logarithmic derivative at $t=1$ gives
\[
\mathrm{avm}(S(T))=\sum_{j=1}^{n-1}\frac{\mu_j}{1+\mu_j}=n-\tr(I+L_T)^{-1}
\]
(the eigenvalue $0$ of $L_T$ contributes $1$ to the trace); this is the case
of a forest of \cite[Lemma~5.1]{AMW}. Finally $S(P_n)=P_{2n-1}$, and $S(T)$
is a path only if $T$ is one. Hence $\tr(I+L_T)^{-1}>\tr(I+L_{P_n})^{-1}$ for
every tree $T\ne P_n$ of order $n$, which is Theorem~\ref{thm:T10}.
(The identity and the unique maximality of the path were also checked
exactly on all trees of orders $2$ to $9$ and $2$ to $16$ respectively;
Section~\ref{sec:comp}.)
\end{remark}

\begin{lemma}[the value for the path]\label{lem:S46}
For $n\ge1$, with $r=\varphi^{-2}$,
\[
\tr(I+L_{P_n})^{-1}=\frac{n}{\sqrt5}+\frac25+\frac{2n}{\sqrt5}\,\frac{r^{2n}}{1-r^{2n}}>\frac{n}{\sqrt5}+\frac25 .
\]
\end{lemma}

\begin{proof}
By (PS), the eigenvalues of $I+L_{P_n}$ are $3-2\cos\theta_j$,
$\theta_j=j\pi/n$, $0\le j\le n-1$. Since $r^2-3r+1=0$ (equivalently
$\varphi^4-3\varphi^2+1=0$), $r+1/r=3$ and
\[
\frac1{3-2\cos\theta}=\frac{r}{1-2r\cos\theta+r^2}
=\frac1{\sqrt5}\Bigl(1+2\sum_{m\ge1}r^m\cos m\theta\Bigr),
\]
by the Poisson kernel, with $r/(1-r^2)=1/(\varphi^2-\varphi^{-2})=1/\sqrt5$;
the series converges absolutely, as $0<r<1$. Now
$\sum_{j=0}^{n-1}\cos(m\theta_j)=\operatorname{Re}\sum_{j=0}^{n-1}\omega^j$ with
$\omega=e^{i\pi m/n}$, and $\omega^n=e^{i\pi m}$. If $2n\mid m$ the sum is $n$.
If $m$ is even and $2n\nmid m$, then $\omega\ne1=\omega^n$ and the sum is $0$.
If $m$ is odd, then $m/n$ is not an even integer, so $\omega\ne1$, and
$\omega^n=-1$; the sum is
$\operatorname{Re}[(1-\omega^n)/(1-\omega)]=\operatorname{Re}[2/(1-\omega)]=1$,
because $\operatorname{Re}[1/(1-e^{i\psi})]=1/2$ whenever $e^{i\psi}\ne1$.
This covers $n=1$ (one eigenvalue, equal to $1$). Summing over $j$ and
interchanging the sums (absolute convergence),
\[
\tr(I+L_{P_n})^{-1}=\frac1{\sqrt5}\Bigl[n+2\sum_{m\text{ odd}}r^m+2n\sum_{l\ge1}r^{2nl}\Bigr],
\]
and $2\sum_{m\text{ odd}}r^m=2r/(1-r^2)=2/\sqrt5$, which gives the formula.
The last term is positive.
\end{proof}

\begin{remark}
The same value follows from \cite[Lemma~4.2, eq.~(8)]{AMW} applied to the
path with $2n-1$ vertices, as the literature audit noted (not checked by a
reviewer).
\end{remark}

\section{Proof of the upper bound}\label{sec:proof}

\begin{theorem}[the upper half]\label{thm:upper}
For every tree $T$ of order $n\ge3$, every maximizing index $p$ of the
normalized coefficients $d_0,\dots,d_{n-2}$ satisfies $p<n-n/\sqrt5+3/5$. In
particular $p\le\lceil n-n/\sqrt5\rceil$; more precisely
$p\le K:=\lceil n-n/\sqrt5-2/5\rceil$.
\end{theorem}

\begin{proof}
Let $\xi=n-n/\sqrt5$, which is irrational since $\sqrt5$ is, and
$K=\lceil \xi-2/5\rceil$; note $K\le\lceil \xi\rceil$, and $K\ge2$ because
$\xi-2/5\ge3(1-1/\sqrt5)-2/5>1$ for $n\ge3$. By Lemma~\ref{lem:S42},
Theorem~\ref{thm:T10} and Lemma~\ref{lem:S46},
\[
\mu_q\le n-\Tr(T)\le n-\Tr(P_n)<\xi-2/5 .
\]
By (D), applied to $q$ (positive coefficients, only real zeros), every
maximizing index $m_q$ of $q_0,\dots,q_{n-1}$ satisfies
$m_q<\mu_q+1<(\xi-2/5)+1$. As $m_q$ is an integer and $\xi-2/5$ is not,
$m_q\le\lceil \xi-2/5\rceil=K$.

If $K\ge n-2$, every maximizing index $p$ of $d_0,\dots,d_{n-2}$ satisfies
$p\le n-2\le K$. Otherwise $0\le K\le n-3$, and $q_{K+1}<q_K$: if not, then,
the ratios $\varrho_j=q_{j+1}/q_j$ being strictly decreasing (proof of
Lemma~\ref{lem:S41}), $\varrho_j>\varrho_K\ge1$ for $j<K$, so
$q_0<q_1<\dots<q_K\le q_{K+1}$, and the maximum of $q_0,\dots,q_{n-1}$ would
be attained at some index $\ge K+1$, contradicting $m_q\le K$ for every
maximizing index $m_q$. Lemma~\ref{lem:S41} with $k=K$ then gives $p\le K$.
In both cases $p\le K\le\lceil \xi\rceil$ and $p\le K<(\xi-2/5)+1=\xi+3/5$.
\end{proof}

\begin{proof}[Proof of the Main Theorem]
Let $T$ be a tree of order $n\ge3$ and $p$ a maximizing index. By
Theorem~\ref{thm:lower}(c), $p\ge\lfloor n/2\rfloor$; by
Theorem~\ref{thm:upper}, $p\le\lceil n-n/\sqrt5-2/5\rceil\le\lceil n-n/\sqrt5\rceil$.
The lower bound is attained by the star, whose unique maximizing index is
$\lfloor n/2\rfloor$ (Lemma~\ref{lem:starshape}). By
Corollary~\ref{cor:pathceil}, for infinitely many $n$ the unique maximizing
index of $P_n$ is $\lceil\xi\rceil$, $\xi=n-n/\sqrt5$, and for those $n$ the
fractional part of $\xi$ exceeds $4/5$, so that
$\lceil\xi-2/5\rceil=\lceil\xi\rceil$: both forms of the upper bound are
attained there.
\end{proof}

\begin{corollary}\label{cor:sharper}
(a) For every tree of order $n\ge3$, every maximizing index is at most
$K=\lceil n-n/\sqrt5-2/5\rceil$, which equals $\lfloor n-n/\sqrt5\rfloor$
when the fractional part of $n-n/\sqrt5$ is below $2/5$. In particular, the
peak of the path can equal $\lceil n-n/\sqrt5\rceil$ only when that
fractional part exceeds $2/5$ (in fact only when it exceeds $4/5$, by
Corollary~\ref{cor:pathceil}).

(b) For an individual tree $T$ of order $n\ge3$, every maximizing index $p$
satisfies $p<1+\sum_j\mu_j/(1+\mu_j)=n+1-\tr(I+L_T)^{-1}$.
\end{corollary}

\begin{proof}
(a) The first statement is Theorem~\ref{thm:upper}. If the fractional part
$f$ of $\xi=n-n/\sqrt5$ is below $2/5$, then
$\xi-2/5\in(\lfloor \xi\rfloor-1,\lfloor \xi\rfloor)$, so $K=\lfloor \xi\rfloor$, which
is less than $\lceil \xi\rceil$ as $\xi$ is irrational. The last sentence follows
since the path is a tree, with Corollary~\ref{cor:pathceil} for the refinement.

(b) Let $k^*$ be the largest maximizing index of $q_0,\dots,q_{n-1}$. By (D)
and Lemma~\ref{lem:S42}, $k^*<\mu_q+1\le n+1-\Tr(T)$. If $k^*\le n-3$, then
$q_{k^*+1}<q_{k^*}$ (as $k^*+1$ is not maximizing), and Lemma~\ref{lem:S41}
with $k=k^*$ gives $p\le k^*$. Otherwise $p\le n-2\le k^*$. In both cases
$p\le k^*<n+1-\Tr(T)$.
\end{proof}

\begin{remark}
For the bound $p\le\lceil n-n/\sqrt5\rceil$ of the conjecture alone,
$\Tr(T)>n/\sqrt5$ is enough: it gives $\mu_q<\xi$, hence $m_q<\xi+1$, so
$m_q\le\lceil \xi\rceil$, and then Lemma~\ref{lem:S41} as in the proof of
Theorem~\ref{thm:upper} with $K$ replaced by $\lceil \xi\rceil$. The extra $2/5$
of Lemma~\ref{lem:S46} is used only for the sharper bound $K$.
Theorem~\ref{thm:T10} is stronger than either requirement. The statement
that the path has the largest peak index among all trees of order $n$ is
not proved here for general $n$, and Theorem~\ref{thm:upper} does not need it
(but see the end of Section~\ref{sec:path}).
\end{remark}

\section{The peak of the path}\label{sec:path}

In this section $n\ge3$, $d_j=d_j(P_n)$, $\xi=n-n/\sqrt5$ and $\eta=n/\sqrt5$
(both irrational). All statements of this section rest on Collins's formula
(C). For every integer $k$ put
\[
h_n(k)=5k^3+2k^2-4k-2-kn^2,
\]
and let $k_0(n)$ be the least integer $k\ge2$ with $h_n(k)>0$.

\begin{lemma}[the peak of the path]\label{lem:S47}
Let $n\ge3$. Then $k_0(n)$ exists and $k_0(n)\le n-1$, and
$d_0,\dots,d_{n-2}$ is strictly increasing up to the index $n-k_0(n)$ and
strictly decreasing after it; so the path has the unique maximizing index
$p(P_n)=n-k_0(n)$. Moreover $p(P_n)\in\{\lfloor \xi\rfloor,\lceil \xi\rceil\}$, and
$p(P_n)=\lceil \xi\rceil$ if and only if $h_n(\lfloor \eta\rfloor)>0$.
\end{lemma}

\begin{proof}
For $2\le k\le n-1$, by (C) (applied to $k$ and to $k+1\le n$) and
$\binom{n+k}{2k+1}/\binom{n+k-1}{2k-1}=\frac{(n+k)(n-k)}{2k(2k+1)}$,
\[
\frac{d_{n-k}}{d_{n-k-1}}=\frac{k-1}{k}\cdot\frac{k+1}{k}\cdot\frac{2k(2k+1)}{(n+k)(n-k)}
=\frac{2(k^2-1)(2k+1)}{k(n^2-k^2)} .
\]
All factors are positive, and $2(k^2-1)(2k+1)-k(n^2-k^2)=h_n(k)$; so
$d_{n-k}>d_{n-k-1}$ if and only if $h_n(k)>0$, and likewise with $<$ and $=$.
Now $h_n(k)/k=5k^2+2k-4-2/k-n^2$ is strictly increasing in real $k>0$, and
$h_n(k)\ne0$ for every integer $k\ge2$: a zero would make $k$ divide $2$
(as $h_n(k)\equiv-2\pmod k$), so $k=2$, and $h_n(2)=38-2n^2\ne0$. Hence,
among integers $k\ge2$, $h_n(k)>0$ exactly for $k\ge k_0(n)$, once $k_0(n)$
exists. And $k_0(n)\le n-1$:
$h_n(n-1)=(n-1)\bigl(5(n-1)^2-n^2\bigr)+2(n-1)^2-4(n-1)-2$ is positive for
$n\ge4$ (both parts are, since $\sqrt5(n-1)>n$, and $2k^2-4k-2>0$ for
$k=n-1\ge3$), and equals $20$ for $n=3$. Reading $k=n-1,n-2,\dots,2$, that is
the index pairs $(n-k-1,n-k)$ from $(0,1)$ upwards, $d$ increases strictly
while $k\ge k_0$ and decreases strictly once $k<k_0$; the unique peak is at
$n-k_0$.

\emph{Location for $n\ge6$.} Then $\eta\ge6/\sqrt5>1+\sqrt2$ (since
$\sqrt5(1+\sqrt2)<5.4$). For $k=\lceil \eta\rceil>\eta$, monotonicity gives
$h_n(k)/k>5\eta^2+2\eta-4-2/\eta-n^2=2\eta-4-2/\eta>0$ (as $\eta^2-2\eta-1>0$), so
$k_0\le\lceil \eta\rceil$. If $k=\lfloor \eta\rfloor-1\ge2$, then $k<\eta-1$ and
$h_n(k)/k<5(\eta-1)^2+2(\eta-1)-4-2/(\eta-1)-n^2=-8\eta-1-2/(\eta-1)<0$, so $h_n<0$ at all
integers $2\le k\le\lfloor \eta\rfloor-1$ and $k_0\ge\lfloor \eta\rfloor$; this is
trivially true also when $\lfloor \eta\rfloor-1<2$. Hence
$k_0\in\{\lfloor \eta\rfloor,\lceil \eta\rceil\}$ and
$p=n-k_0\in\{\lfloor \xi\rfloor,\lceil \xi\rceil\}$, because
$n-\lceil \eta\rceil=\lfloor \xi\rfloor$ and $n-\lfloor \eta\rfloor=\lceil \xi\rceil$.
Since $\lfloor \eta\rfloor\ge2$, $p=\lceil \xi\rceil$ if and only if
$k_0=\lfloor \eta\rfloor$, if and only if $h_n(\lfloor \eta\rfloor)>0$.

\emph{The cases $n=3,4,5$, directly.} Here $\xi\approx1.66,\,2.21,\,2.76$ and
$\lfloor \eta\rfloor=1,1,2$. One has $h_3(2)=20$, $h_4(2)=6$, $h_5(2)=-12$,
$h_5(3)=64$, so $k_0=2,2,3$ and $p=1,2,2=\lfloor \xi\rfloor$. For the criterion
(the polynomial $h_n$ being defined at every integer):
$h_3(1)=-8$, $h_4(1)=-15$, $h_5(2)=-12$ are negative and $p=\lfloor \xi\rfloor\ne\lceil \xi\rceil$,
so ``$p=\lceil \xi\rceil$ if and only if $h_n(\lfloor \eta\rfloor)>0$'' holds, both
sides being false.
\end{proof}

\begin{corollary}[when the path reaches the ceiling]\label{cor:pathceil}
Let $n\ge3$ and let $\theta$ be the fractional part of $\eta=n/\sqrt5$.
(a) If $\theta\ge1/5$, then $p(P_n)=\lfloor \xi\rfloor$. (b) If $\theta<1/10$ and
$n\ge30$, then $p(P_n)=\lceil \xi\rceil$. (c) In particular,
$p(P_n)=\lceil n-n/\sqrt5\rceil$ implies that the fractional part of
$n-n/\sqrt5$ exceeds $4/5$, and $p(P_n)=\lceil n-n/\sqrt5\rceil$ holds for
infinitely many $n$. Thus the upper bound $\lceil n-n/\sqrt5\rceil$ of
Conjecture~\ref{conj:CS} is attained, by the path, for infinitely many $n$,
and for those $n$ it cannot be replaced by $\lfloor n-n/\sqrt5\rfloor$.
\end{corollary}

\begin{proof}
For $n\le5$, Lemma~\ref{lem:S47} gives $p=\lfloor \xi\rfloor$ directly, so (a)
holds and (b), (c) need only $n\ge6$. Let $n\ge6$, so that
$k=\lfloor \eta\rfloor=\eta-\theta\ge2$. Then
\[
h_n(k)/k=2\eta(1-5\theta)+5\theta^2-2\theta-4-2/k .
\]
(a) If $\theta\ge1/5$, the first term is $\le0$, $5\theta^2-2\theta-4<0$ on
$[0,1)$, and $-2/k<0$; so $h_n(\lfloor \eta\rfloor)<0$ and Lemma~\ref{lem:S47}
gives $p=\lfloor \xi\rfloor$. (b) If $\theta<1/10$ and $n\ge30$, the first term
exceeds $\eta\ge30/\sqrt5>13$, while $|5\theta^2-2\theta-4-2/k|<6$; so
$h_n(\lfloor \eta\rfloor)>0$ and $p=\lceil \xi\rceil$. (c) The fractional part of
$\xi=n-\eta$ is $1-\theta$; if $p=\lceil \xi\rceil$ then $\theta<1/5$ by (a), so
$1-\theta>4/5$. Since $1/\sqrt5$ is irrational, (Kr) gives infinitely many
$n$ with $\theta<1/10$, hence infinitely many such $n\ge30$, and for them
$p(P_n)=\lceil \xi\rceil$ by (b). As $p(P_n)$ is the unique maximizing index
(Lemma~\ref{lem:S47}) and $\lceil \xi\rceil>\lfloor \xi\rfloor$, the ceiling cannot
be replaced by the floor for those $n$.
\end{proof}

\emph{Consistency with Theorem~\ref{thm:upper}.} Theorem~\ref{thm:upper} gives
$p(T)\le\lceil \xi-2/5\rceil$ for every tree. When the path is at $\lceil \xi\rceil$,
the fractional part of $\xi$ exceeds $4/5>2/5$, so
$\lceil \xi-2/5\rceil=\lceil \xi\rceil$: no contradiction. When the fractional
part of $\xi$ is below $2/5$ (which implies $\theta=1-\frc{\xi}>3/5\ge1/5$),
Theorem~\ref{thm:upper} gives $p(T)\le\lfloor \xi\rfloor$ for every tree $T$ of
order $n$, and Corollary~\ref{cor:pathceil}(a) gives $p(P_n)=\lfloor \xi\rfloor$; so
for those $n$ the path has the largest peak index among trees of order $n$.
\part*{Discussion, evidence and provenance}

\section{Further remarks and open problems}\label{sec:open}

\begin{enumerate}
\item \emph{Is the path peak-maximal?} Coefficientwise, the star is minimal
(Theorem~\ref{thm:lower}(b)); the project also holds a separate, once
AI-reviewed candidate proof that the path is coefficientwise maximal,
$d_k(T)\le d_k(P_n)$ (not part of this manuscript and not used). Neither
statement locates a peak. Whether every tree of order $n$ peaks no later than
$P_n$ is open in general. By Theorem~\ref{thm:upper} it holds for every $n$
with $p(P_n)=\lceil\xi-2/5\rceil$, $\xi=n-n/\sqrt5$: by
Corollary~\ref{cor:pathceil}(a), for every $n$ for which the fractional part
of $\xi$ is below $2/5$, and also for every $n$ with $p(P_n)=\lceil\xi\rceil$
(for instance $n=9$: $\xi\approx4.975$, $\lceil\xi-2/5\rceil=5$, and
$h_9(4)=10>0$, so $p(P_9)=5$). It remains open for the $n$ with
$p(P_n)<\lceil\xi-2/5\rceil$, that is, $p(P_n)=\lfloor\xi\rfloor$ while the
fractional part of $\xi$ exceeds $2/5$: $806$ of the $n$ with
$3\le n\le2000$ (exact count), the smallest being $3,5,7,8,10$; for $n\le23$
it holds by the exhaustive computation of Section~\ref{sec:comp}. A
mode--mean inequality for $q$ sharper than Darroch's would be one way to
treat these $n$; it is not known to be necessary.
\item \emph{Ratio comparison.} Exact computation supports
$d_{k+1}(T)/d_k(T)\le d_{k+1}(P_n)/d_k(P_n)$ for
$0\le k\le\min(p(P_n),n-3)$ on all trees of order $4$ to $20$ and on $10{,}000$
random trees of orders $24$ to $90$; it is proved here only for $k=0$
($d_1=2n^2-6-\sum_vk_v^2$, which is in the proof of \cite[Theorem~3.3]{AAB}).
The analogous lower comparison with the star,
$d_{k+1}(T)/d_k(T)\ge d_{k+1}(S_n)/d_k(S_n)$ at $k=\lfloor n/2\rfloor-1$, holds
for all trees of order $4$ to $19$ and is false at order $20$: for the double
star $S(9,9)$ at $k=9$ the ratio is $10791/10255<99/94$, the ratio of the star.
\item \emph{Weighted versions.} For a tree $K$ on $j$ vertices with degrees
$k_v$ and a weight $\nu\colon V(K)\to\{0,1,2,\dots\}$, put $|\nu|=\sum_v\nu_v$,
$M=M(K,\nu)$ and $\Phi(K,\nu)$ as in Section~\ref{sec:minors},
$p^\nu=\one+x\tau-x\nu$ (so $p^\nu_v=1+(2-k_v-\nu_v)x$), and
\[
R(K,\nu)=(|\nu|+j-1)(1+x)\,\Phi(K,\nu)-p^{\nu\T}\adj(M)\,p^\nu ;
\]
for $\nu=we_u$ this is $R_w(K,u)$ of Section~\ref{sec:rooted}.
Theorem~\ref{thm:T5}(iii) with Remark~\ref{rem:Ccone} is the case
$\nu=e_u$ of the statement ``$R(K,\nu)\in\mathcal C_{j+1}$ for every
weighted tree with $\nu\ne0$'', which is open for general weights; some
stronger auxiliary statements of this kind are false (the project records
counterexamples on stars with heavy leaves).
\item \emph{Margins.} For even $n$ the relative margin $1-d_{m-1}/d_m$ tends
to $0$ (it is of order $1/n$ for double stars and much smaller, about
$10^{-6}$ at $n=10^5$, for balanced caterpillars with a slowly growing spine);
the cone argument nevertheless gives a uniform proof.
\end{enumerate}

\section{Computational evidence}\label{sec:comp}

The computations below support, but are not part of, the proofs. All are
in exact arithmetic (integers and rationals) unless stated. Paths are
relative to the repository of the project.

\begin{enumerate}
\item \emph{The conjecture itself.} Every maximizing index lies in
$[\lfloor n/2\rfloor,\lceil n-n/\sqrt5\rceil]$ for all $23{,}942{,}357$ trees of
order $3$ to $23$ (\nolinkurl{experiments/CJ-0001-collins-shor-peak/exhaustive.py});
orders $21$ to $23$, beyond the range $n\le20$ reported in~\cite{AAB}, were
computed by one implementation only. One tie occurs (order $13$, indices $6$
and $7$, common value $d_6=d_7=189972$; the tie and its value were found
again by the independent reviewer of revision~2 and by a separate
recomputation over all $1{,}301$ trees of order $13$); the peak is always
$\lfloor n/2\rfloor$ or $\lfloor n/2\rfloor+1$ in this range, and in every
order the largest maximizing index over all trees is that of the path.
\item \emph{Lower half, notes' code} (\nolinkurl{lower_bound.py}, result
\nolinkurl{results/lower-bound.json}, run on a clean tree at commit
\texttt{b1359cc}): $1{,}040{,}693$ checks of the identities and memberships of
Sections~\ref{sec:distances} and~\ref{sec:lowerproof} on all rooted trees of
order at most $9$ and all merges of order at most $9$; star minimality
$\ell_T-\ell_{S_n}\in\mathcal C_{n-1}$ for all $1{,}346{,}022$ trees of order
$3$ to $20$ (\nolinkurl{leaf_telescoping.py}; the project's notes give
$1{,}346{,}023$, which is the number of trees of order $2$ to $20$). No
failure.
\item \emph{Lower half, independent check of this text}
(\nolinkurl{manuscript_check.py}, written for this manuscript from its
definitions, without the notes' code): Lemma~\ref{lem:star} for $w=0,1,2$,
Proposition~\ref{prop:rules}, Corollary~\ref{cor:leaf},
Propositions~\ref{prop:S33}--\ref{prop:S35}, the memberships of
Theorem~\ref{thm:T5} and Corollary~\ref{cor:phi1} on all rooted trees of
order at most $8$ and all merges of order at most $8$, and
Theorem~\ref{thm:lower} on all trees of order $2$ to $12$: $11{,}379$ checks,
no failure. The result file records the command and revision.
\item \emph{Upper half} (\nolinkurl{upper_bound.py}, result
\nolinkurl{results/upper-bound.json}, run at commit \texttt{917d657} on a
clean worktree, 0 failures): Theorem~\ref{thm:T10} on all trees of order $2$
to $20$; the inequality of Lemma~\ref{lem:S46} exactly for $n\le200$;
Theorem~\ref{thm:upper} on all trees of order $3$ to $18$ and, through (C), on
paths for $n\le2000$; Lemma~\ref{lem:S45} on all rooted trees with
$2\le|X|\le8$, $j_1+j_2\le8$; Lemma~\ref{lem:S44}(iii) for $N\le60$;
Lemma~\ref{lem:S43} on all pairs of rooted trees of order at most $6$;
Lemma~\ref{lem:S42} in floating point (exploratory) on trees of order $3$ to
$14$. The criterion of Lemma~\ref{lem:S47} was checked against the exact peak
of the path for $n\le2000$.
\item \emph{Revision~3} (\nolinkurl{amw_reduction_check.py}, result
\nolinkurl{results/amw-reduction-check.json}, exact): for
Remark~\ref{rem:AMW}, the identities $m_k(S(T))=[t^k]\det(I+tL_T)$ and
$\mathrm{avm}(S(T))=n-\tr(I+L_T)^{-1}$ on all $94$ trees of order $2$ to $9$,
the path as the unique maximizer of $\mathrm{avm}$ on all $32{,}507$ trees of
order $2$ to $16$, and $\mathrm{avm}(P_N)>aN+b$ for $4\le N\le200$ (exact
sign test in $\mathbb Q(\sqrt5)$); the count of open orders in
Section~\ref{sec:open}, item~1 (integer arithmetic, from
Lemma~\ref{lem:S47}); and the tie at order~$13$. In addition, eq.~(8)
of~\cite{AMW} was compared with the exact value of $\mathrm{avm}(P_N)$ for
$N\le80$ in $60$-digit floating point (exploratory; agreement to $50$
digits).
\end{enumerate}
Finite checks establish the statements only in the range checked.

\section{Formal verification}\label{sec:formal}

Only auxiliary statements are formalized in Lean~4 with Mathlib (directory
\nolinkurl{formalization/Conjectura/Research/CollinsShorPeak/}): a lemma on
the maximizing indices of positive log-concave sequences (used in the first
version of this text, no longer used); the cones $\K_d$ with (K1), (K2) and
Lemma~\ref{lem:A}; the cone $\mathcal C_d$ with its product rules; and the
sequence part of Lemma~\ref{lem:S41} (positivity of $q$ and the decrease of
its ratios as hypotheses). Nothing about trees, determinants, distance
matrices, Darroch's rule or the trace of $(I+L)^{-1}$ is formalized. The
project's receipts record that the library builds without \texttt{sorry} and
with only the standard axioms; the comparison of the Lean statements with the
written ones has been made by the agents that wrote them, not independently.

\section{Literature and novelty}\label{sec:literature}

\emph{Coverage.} The project's searches (November 2026, by AI literature
auditors) used general web search and arXiv only (the queries are listed verbatim in
the project's dossier for this problem), the papers citing~\cite{AAB} listed by Semantic Scholar and
OpenAlex (five papers: arXiv:2407.03309, 2206.07561, 2104.14516, 2103.00647,
1812.05734; none resolves the conjecture), \cite{AAB} and~\cite{Col89} read in full, Collins's
thesis~\cite{Col86} in part, \cite{GL78} in part, \cite{Mer90}, \cite{AMW}
(arXiv version), \cite{Pit97} (one page), \cite{CS06}, and the papers of
Mohar and Csikv\'ari on Laplacian coefficients of trees; for revision~3, the
$14$ papers citing~\cite{AMW} and the $63$ works since 2023 citing~\cite{GL78}
listed by OpenAlex (by title), arXiv searches for 2023--2026, \cite{GP71} and
p.~83 of~\cite{GL78} in scans. Related recent work, not used here: a
rooted-spanning-forest formula for the coefficients of the distance
polynomial of a tree (Richman, Shokrieh and Wu, arXiv:2411.11488,
Theorem~C), and Chebotarev's study of the forest matrix of the path and the
golden ratio (Discrete Appl.\ Math.\ 156 (2008) 813--821; abstract only),
which may contain the values of Lemmas~\ref{lem:S44} and~\ref{lem:S46}.
\emph{Not searched:} MathSciNet, zbMATH, Google Scholar, MathOverflow; the
full texts of the papers citing~\cite{AMW}.
An unsuccessful search does not establish originality, and nothing below is
a claim of originality.

\smallskip
\noindent\textbf{Known results used or reproved here.}
\begin{itemize}
\item Lemma~\ref{lem:ident}, (GL), (GP), (LC): \cite{GL78,GP71,AAB}.
Lemma~\ref{lem:L1} is an immediate consequence of (GL) and (GP).
\item Lemma~\ref{lem:path}, first equality: a consequence of the path
expansion of cofactors of Maybee, Olesky, van den Driessche and
Wiener~\cite{MOvdDW}, as quoted in~\cite[eq.~(4)]{CS06} (the specialization to
trees was made by the project's auditor and is unreviewed). The membership
part rests on Lemma~\ref{lem:K}, for which no source was searched.
\item Corollary~\ref{cor:leaf}(d) and the peak of the star: \cite[Theorem~1]{Col89}.
A leaf recursion for the minors of $D-\lambda I$ is
\cite[Theorem~4]{Col86}; the form of Corollary~\ref{cor:leaf}(b), with
$1+x$ and a polynomial $r_u$ with non-negative coefficients, was not found
there.
\item Lemma~\ref{lem:S40} and the interlacing behind Lemma~\ref{lem:S42}:
Merris~\cite{Mer90}. Darroch's rule: \cite{Dar64,Pit97}.
\item Theorem~\ref{thm:T10}: a consequence of \cite[Theorem~4.4 with
Lemma~4.2]{AMW} through the subdivision of the tree (Remark~\ref{rem:AMW},
where the reduction is proved; the reduction was first derived by the
project's literature auditor and confirmed by the reviewer of revision~2).
Our proof by merging pendant paths is independent of theirs; as a statement,
Theorem~\ref{thm:T10} is a rediscovery. Lemma~\ref{lem:S46} also follows
from \cite[Lemma~4.2]{AMW}.
\item The path: (C) is~\cite[Theorem~3]{Col89} (numbering of the journal
version, read in the project; an independent reviewer reports that the
author's preprint numbers it Theorem~2), whose location of the peak up to one
has an informal last step; Lemma~\ref{lem:S47} makes it exact.
\end{itemize}

\noindent\textbf{No prior source found within the coverage.}
Both halves of Conjecture~\ref{conj:CS} (beyond the bound
$\lceil2n/3\rceil$ of~\cite{AAB} and the families treated in~\cite{Col89});
the coefficientwise minimality of the star, Theorem~\ref{thm:lower}(b) (the
classical extremal results for Laplacian coefficients of trees concern
$\det(I+xL)$ alone and do not give it); the rooted calculus of
Section~\ref{sec:rooted}, Theorem~\ref{thm:T5} and the sum of squares of
Proposition~\ref{prop:S35}; the use of Darroch's rule for the peak of the
distance polynomial; the exact criterion of Lemma~\ref{lem:S47} and
Corollary~\ref{cor:pathceil} (no separate search was made for these two).

\noindent\textbf{Open auxiliary conjectures} are those of
Section~\ref{sec:open}, items 1 to 3; none is used in the proofs.

\section{Status of the claims}\label{sec:status}

\begin{itemize}
\item \emph{Mathematical status:} candidate proof. Each half has a complete
written argument; it is not an established result.
\item \emph{Reviews of the sources of this text} (all by AI models of one
vendor, in fresh contexts that received the claim and the source but not the
author's reasoning): the rooted calculus of
Section~\ref{sec:rooted} in its original weighted form (one review, model
\nolinkurl{claude-opus-5-5}); Lemma~\ref{lem:K} (one complete audit,
\nolinkurl{claude-opus-5-5}); Lemma~\ref{lem:path} (without weights: one
review and one audit, \nolinkurl{claude-opus-5-5}; with weights: within the
reviews of Sections~\ref{sec:distances} and~\ref{sec:lowerproof}); Sections~\ref{sec:distances}
and~\ref{sec:lowerproof} (two reviews, \nolinkurl{claude-opus-5-5} and
\nolinkurl{claude-fable-5-2}, and a recheck of the revised text by
\nolinkurl{claude-sonnet-5-5}); Sections~\ref{sec:spectral}--\ref{sec:proof}
(two reviews, \nolinkurl{claude-opus-5-5} and \nolinkurl{claude-sonnet-5-5},
and a recheck); Section~\ref{sec:path} (one review,
\nolinkurl{claude-sonnet-5-5}). All reported no gap and listed defects of
wording, citation or bookkeeping, which are corrected here.
\item \emph{What is new in this text and not reviewed:} the presentation of
Sections~\ref{sec:minors}--\ref{sec:rooted} with root weights only, in which
the rules for raising a weight and planting are derived from
Lemma~\ref{lem:star} (in the project's notes they came from separate
recurrences for weighted trees), and the derivation of the leaf recursion
from the plant rule; Lemma~\ref{lem:star} itself for root weights; the
assembly of the two halves. These were checked by exact computation
(Section~\ref{sec:comp}, item 3) and, in the first version of this text
(commit \texttt{8df6bea}), by two independent AI reviews of the whole
document (models \nolinkurl{claude-opus-5-5} and
\nolinkurl{claude-sonnet-5-5}, fresh contexts, same vendor as the author;
both ``no gap found within stated scope'', with two errors in remarks or
proofs not used by the Main Theorem and several defects of wording and
bookkeeping, corrected in revision~2).
\item \emph{Revision~2} (commit \texttt{87c3c8c}) added
Lemma~\ref{lem:starshape} (the shape of the star's sequence), which replaces
the use of (LC) and of \cite[Theorem~1]{Col89} in
Theorem~\ref{thm:lower}(c) and in the attainment of the lower bound, and
removed the lemma on log-concave sequences; Lemma~\ref{lem:starshape} was
checked exactly for $3\le n<800$. Revision~2 was rechecked
(\nolinkurl{claude-sonnet-5-5}, its changes only) and then reviewed as a
whole by an AI model of another vendor (OpenAI, through Codex; the project
director's description of the model is ``GPT Astra 6 (Extra High)'', the
report's own is ``agent based on GPT-6''), which reconstructed Sections
2--11, ran its own exact checks, and found no gap; it reported one error in a
remark not used by the Main Theorem (the multiplicity of the eigenvalue $0$
in the spectrum of the subdivision, Remark~\ref{rem:AMW}) and five defects
of wording or citation.
\item \emph{Revision~3} (this text) corrects all six: Remark~\ref{rem:AMW}
now proves that Theorem~\ref{thm:T10} follows from~\cite{AMW}; $\varphi$ is
defined; the Graham--Lov\'asz peak is written with the floor; item~1 of
Section~\ref{sec:open} states exactly which orders remain open and no longer
claims that a sharper mode--mean inequality is needed; $R(K,\nu)$ is defined
in item~3 there; and the decimal estimate at the end of the proof of
Lemma~\ref{lem:S44} is replaced by an exact one. No step used by the Main
Theorem changed. The recheck of revision~3 is recorded in the project's
journal, not here.
\item \emph{Human review:} none. \emph{Formal verification:} partial, of
auxiliary lemmas only (Section~\ref{sec:formal}). \emph{Novelty:} as in
Section~\ref{sec:literature}.
\end{itemize}

\section*{Use of AI systems}

\emph{Systems.} Claude (Anthropic): models \nolinkurl{claude-fable-5-2},
\nolinkurl{claude-opus-5-5} and \nolinkurl{claude-sonnet-5-5}, run through
Claude Code in the Conjectura repository, 26--27 November 2026.

\emph{What they did.} \nolinkurl{claude-fable-5-2}: the reduction of the lower
half to rooted trees, the rooted calculus in its weighted form, the leaf
recursion, Theorem~\ref{thm:T5} and Theorem~\ref{thm:lower} (session of
26 November 2026), the exhaustive computations to order $23$, and parts of
the literature audits. \nolinkurl{claude-opus-5-5}: the spectral route to the
upper half (Lemmas~\ref{lem:S40}--\ref{lem:S47}, Theorems~\ref{thm:upper}
and~\ref{thm:T10}), the stored computations for it, the formalization, the
reviews named in Section~\ref{sec:status}, and this text, including the
restatement of the lower half with root weights. \nolinkurl{claude-sonnet-5-5}:
reviews, rechecks and computations. An AI model of OpenAI, run by the
project director through Codex, reviewed revision~2 independently (its
report is kept in the project's records). No human contributed mathematics.

\emph{Human involvement.} The project director chose the problem and set the
objectives, budgets and policies of the sessions. No mathematician has read
the proofs.

\section*{Attribution}

Research direction: Gustavo Bernardo Brito (UFG), director of the Conjectura
project. Mathematical content: AI systems as listed above. External
mathematical results: those cited, in particular
\cite{GL78,GP71,AAB,Col89,Dar64,Pit97,Mer90,AMW}. Authorship of any eventual
publication is to be decided by the director under the project's
attribution policy; operating an AI system does not by itself confer
authorship, and AI systems are not listed as authors.

\section*{Limitations}

\begin{itemize}
\item The proofs are candidate proofs, checked by AI reviews of their parts
and by exact computation; the first version of this document was reviewed
as a whole by two AI models of one vendor, and revision~2 by one AI model of
another vendor (Section~\ref{sec:status}); no human has read it. AI review
is not peer review, and agreement between AI reviewers does not establish
correctness.
\item Several classical results are cited without proof, and some were not
read at their sources (Section~\ref{sec:tools}): Darroch's paper and the
textbook facts.
\item Section~\ref{sec:path} rests on Collins's formula (C), checked against
the distance matrix only for $n\le30$.
\item Novelty is not established; Theorem~\ref{thm:T10} is a consequence
of a published theorem (Remark~\ref{rem:AMW}).
\end{itemize}

\appendix

\section{Correspondence with the project's records}\label{app:map}

For reviewers who wish to trace a statement to the project's notes
(directory \nolinkurl{research/active/CJ-0001-collins-shor-peak/notes/}) and to
the reviews of those notes (\nolinkurl{reviews/}):

\smallskip
\noindent\begin{tabular}{@{}p{0.37\textwidth}p{0.58\textwidth}@{}}
Lemma~\ref{lem:L1} & L1 (\nolinkurl{lemmas.md})\\
Lemma~\ref{lem:starshape} & new in revision~2 of this text (S23(d) and Collins's Theorem~1 give the value at the middle and the shape)\\
Lemma~\ref{lem:A} & Lemma A(a) (\nolinkurl{wiener-form.md})\\
Lemma~\ref{lem:K}, Lemma~\ref{lem:path} & Lemma K, Lemma S24 (\nolinkurl{wiener-form.md}); Lemma S32 (\nolinkurl{lower-bound.md})\\
Corollary~\ref{cor:phi1} & proof of T6(a) (\nolinkurl{lower-bound.md})\\
Lemma~\ref{lem:star}, Prop.~\ref{prop:rules} & (\(\star\)) and S19 (\nolinkurl{rooted-calculus.md}), restricted to root weights; S15, S16 (\nolinkurl{gamma-components.md}) are replaced by Lemma~\ref{lem:star}\\
Corollary~\ref{cor:leaf} & S23 (\nolinkurl{leaf-telescoping.md})\\
Props.~\ref{prop:S33}--\ref{prop:S35}, Thm.~\ref{thm:T5} & S33--S35, T5 (\nolinkurl{lower-bound.md}); registry MD-0019\\
Thm.~\ref{thm:lower}, Lemma~\ref{lem:starvalue} & T6 (\nolinkurl{lower-bound.md}), S23(d)\\
Lemmas~\ref{lem:S40}--\ref{lem:S46}, Thms.~\ref{thm:upper}, \ref{thm:T10} & S40--S46, T9, T10 (\nolinkurl{upper-bound.md}); MD-0022, MD-0023\\
Lemma~\ref{lem:S47}, Cor.~\ref{cor:pathceil} & S47, S48 (\nolinkurl{path-peak.md}); MD-0024\\
\end{tabular}

\begin{thebibliography}{99}

\bibitem{AAB}
G.~Aalipour, A.~Abiad, Z.~Berikkyzy, L.~Hogben, F.~H.~J. Kenter, J.~C.-H. Lin
and M.~Tait, \emph{Proof of a conjecture of Graham and Lov\'asz concerning
unimodality of coefficients of the distance characteristic polynomial of a
tree}, arXiv:1507.02341; Electron.\ J.\ Linear Algebra 34 (2018) 373--380
(journal details from secondary sources, not checked against the journal).

\bibitem{AMW}
E.~O.~D. Andriantiana, V.~Razanajatovo Misanantenaina and S.~Wagner,
\emph{The average size of matchings in graphs}, Graphs Combin.\ 36 (2020)
539--560; arXiv:1903.09825.

\bibitem{BP}
A.~Berman and R.~J. Plemmons, \emph{Nonnegative Matrices in the
Mathematical Sciences}, SIAM, Philadelphia, 1994.

\bibitem{BH}
A.~E. Brouwer and W.~H. Haemers, \emph{Spectra of Graphs}, Springer, New
York, 2012.

\bibitem{CS06}
P.~Chebotarev and E.~Shamis, \emph{Matrix-forest theorems},
arXiv:math/0602575.

\bibitem{Col86}
K.~L. Collins, \emph{Distance matrices of trees}, Ph.D. thesis,
Massachusetts Institute of Technology, 1986.

\bibitem{Col89}
K.~L. Collins, \emph{On a conjecture of Graham and Lov\'asz about distance
matrices}, Discrete Appl.\ Math.\ 25 (1989) 27--35.

\bibitem{Dar64}
J.~N. Darroch, \emph{On the distribution of the number of successes in
independent trials}, Ann.\ Math.\ Statist.\ 35 (1964) 1317--1321.

\bibitem{GL78}
R.~L. Graham and L.~Lov\'asz, \emph{Distance matrix polynomials of trees},
Adv.\ Math.\ 29 (1978) 60--88.

\bibitem{GP71}
R.~L. Graham and H.~O. Pollak, \emph{On the addressing problem for loop
switching}, Bell System Tech.\ J.\ 50 (1971) 2495--2519.

\bibitem{HLP}
G.~H. Hardy, J.~E. Littlewood and G.~P\'olya, \emph{Inequalities}, 2nd ed.,
Cambridge University Press, 1952.

\bibitem{HJ}
R.~A. Horn and C.~R. Johnson, \emph{Matrix Analysis}, 2nd ed., Cambridge
University Press, 2013.

\bibitem{MOvdDW}
J.~S. Maybee, D.~D. Olesky, P.~van den Driessche and G.~Wiener,
\emph{Matrices, digraphs, and determinants}, SIAM J.\ Matrix Anal.\ Appl.\
10 (1989) 500--519.

\bibitem{Mer90}
R.~Merris, \emph{The distance spectrum of a tree}, J.\ Graph Theory 14
(1990) 365--369.

\bibitem{Pit97}
J.~Pitman, \emph{Probabilistic bounds on the coefficients of polynomials
with only real zeros}, J.\ Combin.\ Theory Ser.\ A 77 (1997) 279--303.

\end{thebibliography}

\end{document}
